\documentclass[12pt,a4paper]{article}
\usepackage{amsmath,amssymb,amsthm}
\usepackage{geometry}
\usepackage{hyperref}
\usepackage{microtype}
\usepackage{booktabs}
\usepackage{array}
\usepackage{enumitem}
\geometry{margin=2.5cm}

\newtheorem{theorem}{Theorem}[section]
\newtheorem{lemma}[theorem]{Lemma}
\newtheorem{corollary}[theorem]{Corollary}
\newtheorem{definition}[theorem]{Definition}
\newtheorem{remark}[theorem]{Remark}
\newtheorem{observation}[theorem]{Observation}

\title{The Cascade Series --- Part II\\[0.5em]
\large Quantum Mechanics from the Cascade:\\
Effective Theory of a 4-Dimensional Observer\\
in the Sphere-Area Geometry}
\author{RTAC}
\date{March 2026}

\begin{document}
\setlength{\emergencystretch}{3em}
\maketitle

\begin{abstract}
The cascade series tests one hypothesis: the infinite-dimensional unit ball, descended to
four dimensions, is indistinguishable from our universe. The companion paper~\cite{paperI}
derived a geometric invariant $I = 1.0990\times 10^{-120}$ from the sphere-area cascade
using only orthogonality as input. Here we show that a 4-dimensional observer embedded in
the cascade geometry recovers the full structural framework of quantum mechanics---including
the Born rule, complex amplitudes, non-commutativity, the Schr\"{o}dinger equation,
entanglement, and Bell inequality violation---without quantum postulates.

The argument proceeds in six steps, each a theorem about spheres and Beta functions.
(1)~The cascade's boundary dominance ($\Omega_{d-1}/V_d = d$) fixes the state space as the
unit sphere $S^{d-1}$; projection onto a 4D subspace produces Gaussian statistics via
concentration of measure. (2)~The cascade's orthogonal slicing structure induces measurement
bases; the commutator of rank-1 projections at angle $\theta$ is
$\|[P,Q]\| = \tfrac{1}{2}|\sin 2\theta|$, vanishing at $\theta\in\{0,\pi/2,\pi\}$.
(3)~The Born rule $p = \cos^2\theta$ is derived from the geometry of the spherical cap, the
slicing integrand $(1-z^2)^{(d-2)/2}$, and Parseval's identity $\sum(u\cdot e_j)^2 = 1$;
no probability axiom enters. (4)~The precession angle $\alpha = \pi/2$ between consecutive
slicing axes is forced by the same orthogonality axiom that forces $\sqrt{\pi}$ in the
slicing recurrence; two quarter-turns give $J^2 = -\mathrm{Id}$, producing complex
amplitudes. The $U(1)$ generated by $J$ is forced as a gauge equivalence under the
cascade hypothesis (Theorem~\ref{thm:U1-gauge-forced}), making the projective state
space $\mathbb{C}P^{n-1}$ via the Hopf bundle and the Born rule the complex
$|\langle u, v\rangle_{\mathbb{C}}|^2$; the cascade is committed to complex-amplitude QM
and accordingly to the Renou \emph{et al.}~(2021)~\cite{renou2021} tripartite Bell
statistics. The phase--obstruction lockstep (Corollary~\ref{cor:lockstep}) follows:
the propagator phase is imaginary if and only if $S^{d-1}$ carries a hairy ball zero,
both controlled by the parity of $d$. The cascade's classification has period~4; the
Clifford refinement to period~8 is established in Part~III.
(5)~Time is dimensional collapse: each cascade slicing operation resolves one dimension
of $B^\infty$, with the cascade's truncation height advancing by one and the observer's
$S^3$-host shell taking one asymptotic step toward the cover-sheet pseudo-horizon.
The discrete propagator
$K = \prod|L(j)|\cdot i^{D-d'}$ is exact; the Schr\"{o}dinger equation $i\dot\psi =
\mathcal{H}\psi$ with $\mathcal{H} = (1-N)/N^2$ (a Gamma function ratio) is the effective
description for a 4D observer who cannot resolve individual cascade steps. Each step
compactifies one direction to $R_{\rm eff} = 1/\sqrt{d+3}$ (from the Beta function),
producing geometric decoherence: the decoherence factor
$D = \exp(-\Delta x^2/4R_{\rm eff}^2)$ identifies the compactification radius as the
coherence length. Over the full 213-step descent, the cascade retains $93.9\%$ coherence;
the $6.1\%$ eigenvalue deficit is the inter-layer coupling that drives the descent
corrections in the Part~0 Supplement. (6)~Iterated slicing gives a tensor product
structure with generic entanglement, and the CHSH correlator computed from the Born rule
on $S^7$ gives $S = 2\sqrt{2}$, violating the classical bound of~2.

The value of $\hbar$ is not derived; it enters as the ratio of the forced precession
scale $\pi/2$ to the physical time increment, encoding only the unit matching between
geometric and physical time. No physical law or free geometric parameter beyond the
observer's dimensionality enters the derivation.
\end{abstract}

\tableofcontents
\newpage

\section{Introduction and Summary of the Foundation}

In~\cite{paperI}, the sphere-area cascade was derived from a single axiom: orthogonality.
Starting from the infinite-dimensional unit ball---which has zero volume, zero surface area,
and no interior---the cascade compresses orthogonal directions one at a time, descending
from $d=\infty$. The unique dimension-independent constant in the slicing recurrence is
$\Gamma(\tfrac{1}{2})=\sqrt{\pi}$, forced by the quarter-turn integral $B(\tfrac{1}{2},\cdot)$.
This constant enters in two roles (multiplicative factor and bare decay rate), producing two
distinguished dimensions $d_1=19$ and $d_2=217$, and a cascade invariant
$I=9\,\Omega_{19}\,\Omega_{217}/\pi^2=1.0990\times 10^{-120}$.

Corollary~3.2 of~\cite{paperI} establishes that sphere areas are the unique independent
cascade quantities: every other cascade object is derived from sphere areas, and any
identification mapping cascade content to energy must do so with a universal constant.
Section~\ref{sec:compactification} of this paper develops the full compactification
interpretation: each slicing step sends one direction's effective radius to
$R_{\rm eff}(d)=1/\sqrt{d+3}$, and the boundary-to-volume ratio satisfies
$\Omega_{d-1}/V_d=d$ for all $d\geq 1$. The cascade's content is its boundaries; sphere
areas, not volumes, are the primary objects.

The present paper asks: what does the cascade geometry look like from inside? Specifically,
if an observer is restricted to a 4-dimensional cross-section of the full cascade structure,
what measurement theory does that observer use? We show the answer is quantum mechanics.

The logical structure is:
\begin{itemize}
\item \emph{Foundation~\cite{paperI}}: Orthogonality $\to$ $\sqrt{\pi}$ $\to$ slicing
recurrence $\to$ cascade $\to$ tower $(d_0,d_1,d_2)$ $\to$ $I\approx 10^{-120}$. Each step
is an asymptotic compactification with $R_{\rm eff}=1/\sqrt{d+3}$.
\item \emph{This paper}: Cascade geometry + 4D observer $\to$ Hilbert space + Gaussian
measure + Born rule ($p=\cos^2\theta$) + complex amplitudes ($J^2=-\mathrm{Id}$) +
discrete propagator + decoherence ($D=e^{-\Delta x^2/4R_{\rm eff}^2}$) + Bell violation
($S=2\sqrt{2}$) $=$ QM.
\end{itemize}

The derivation does not explain why we observe four dimensions. That is an empirical input,
just as in~\cite{paperI}. What we show is that, given a 4D perspective, the cascade's
internal structure reproduces the quantum formalism.

\section{The Observer's Arena}

\subsection{State space is the unit sphere}

Quantum states are unit vectors in a Hilbert space $\mathcal{H}$. The set of pure states is
the unit sphere $S(\mathcal{H})$. In finite dimension $d$, this is $S^{d-1}\subset\mathbb{R}^d$
(or $S^{2d-1}\subset\mathbb{C}^d$ after complexification). The cascade geometry provides
exactly this arena: the unit sphere $S^{d-1}$ at each level $d$ of the descent. A 4D
observer---one who can access only 4 orthogonal directions---sees the projection of
$S^{d-1}$ onto a 4-dimensional subspace. This projected space is the observer's effective
state space.

\subsection{Why the cascade fixes the arena}

Standard quantum mechanics postulates a Hilbert space but does not derive it. The cascade
provides the derivation: the arena is the unit sphere because the cascade is a sequence of
volume-preserving orthogonal compressions of the unit ball. Every state accessible to any
observer at any dimension lies on a unit sphere. The constraint $|\psi|^2=1$, which in QM
is the normalisation postulate, is here a theorem: the cascade operates on $B^d$, whose
boundary is $S^{d-1}$.

\subsection{Boundary dominance and the primacy of the sphere}

The arena is the unit sphere rather than the unit ball for a reason that~\cite{paperI},
Section~3 makes precise.

\begin{theorem}[Boundary dominance; Theorem~3.1 of~\cite{paperI}]
$\Omega_{d-1}/V_d=d$ for all $d\geq 1$.
\end{theorem}

At $d=4$: $\Omega_3/V_4=4$, so the boundary $S^3$ carries a fraction $d/(d+1)=4/5=80\%$
of the content of $B^4$. At the cascade dimensions $d_1=19$ and $d_2=217$, the fraction
on the boundary is $19/20=95\%$ and $217/218\approx 99.5\%$ respectively. The cascade's
content is its boundaries, increasingly so at each step. The observer's state space is the
sphere because the sphere carries nearly all the geometric content available to the observer.
This is not an axiom; it is a consequence of the Gamma function identity
$\Omega_{d-1}/V_d = 2\pi^{d/2}/[\Gamma(d/2)] \div \pi^{d/2}/[\Gamma(d/2+1)] = d$.

\begin{remark}[Layered boundary dominance]
The cascade stacks boundary-dominant steps: the $d_2$-sphere's boundary encloses the
$d_1$-sphere's boundary, which encloses the $d_0$-sphere's boundary, which encloses the
observer's $S^3$. At each step, the boundary $S^{d-1}$ carries a fraction $d/(d+1)$ of the
content. The cascade is a nested sequence of shells, each carrying nearly all the geometric
content of its level, with the observer on the innermost shell.
\end{remark}

\section{Gaussian Statistics from Projection}

\subsection{Concentration of the slicing integrand}

The slicing integrand $f_d(x) = (1-x^2)^{d/2}$ that defines the cascade recurrence
is itself approximately Gaussian. This is the concrete mechanism underlying the
projection theorem that follows.

\begin{lemma}[Gaussian concentration of the integrand]\label{lem:gaussian}
The slicing integrand $f_d(x) = (1-x^2)^{d/2}$ satisfies:
\begin{enumerate}
\item[(a)] $f_d(x) = \exp(-dx^2/2 + O(dx^4))$, approximating a Gaussian with standard
deviation $\sigma(d) = 1/\sqrt{d}$;
\item[(b)] the half-maximum width is $\sqrt{2\ln 2}/\sqrt{d}$;
\item[(c)] as $d\to\infty$, $f_d\to\delta(x)$.
\end{enumerate}
\end{lemma}
\begin{proof}
Expanding: $\ln(1-x^2) = -x^2 - x^4/2 - x^6/3 - \cdots$ for $|x|<1$. Therefore
$(d/2)\ln(1-x^2) = -dx^2/2 + O(dx^4)$. The half-maximum occurs at
$x_{1/2} = \sqrt{1-2^{-2/d}} \approx \sqrt{2\ln 2}/\sqrt{d}$.
As $d\to\infty$, $x_{1/2}\to 0$ and $\int f_d\,dx \approx \sqrt{2\pi/d}\to 0$,
confirming convergence to $\delta(x)$.
\end{proof}

This lemma is the elementary origin of all Gaussian structure in the cascade. The
projection theorem (below) elevates it to a distributional statement on the unit
sphere; the compactification results (Section~\ref{sec:compactification}) use it to
define the effective radius of each temporal step.

\subsection{The projection theorem}

\begin{theorem}[Gaussian emergence]\label{thm:gaussian}
Let $x$ be uniformly distributed on $S^{d-1}(\sqrt{d})\subset\mathbb{R}^d$. Let
$\pi_k:\mathbb{R}^d\to\mathbb{R}^k$ be projection onto any $k$-dimensional subspace, with
$k$ fixed and $d\to\infty$. Then $\pi_k(x)$ converges in distribution to
$\mathcal{N}(0,I_k)$, the standard $k$-dimensional Gaussian.
\end{theorem}

This is a consequence of the Poincar\'{e} limit theorem (see~\cite{ledoux,milman}). The
rate of convergence is $O(1/d)$, so for the cascade dimensions $d_1=19$, $d_2=217$, the
approximation is already excellent for $k=4$.

\subsection{The cascade's Gaussian is the QM wavefunction}

The Gaussian that emerges from projection is structurally identical to the ground-state
wavefunction of the quantum harmonic oscillator:
$\psi_0(x)=\pi^{-1/4}\exp(-x^2/2)$. The normalisation constant $\pi^{-1/4}$ involves
exactly the $\sqrt{\pi}$ from the cascade. In standard QM this constant is obtained by
requiring $\int|\psi|^2dx=1$ with the Gaussian integral $\int\exp(-x^2)dx=\sqrt{\pi}$.
In the cascade, $\sqrt{\pi}$ is not a normalisation convention: it is the unique
dimension-independent constant forced by orthogonal compression (Theorem~3.1
of~\cite{paperI}). The cascade derives what QM postulates.

\subsection{Why Gaussians are universal}

The cascade provides a structural explanation for why Gaussians dominate quantum physics.
Any 4D observer embedded in a high-dimensional unit-sphere geometry necessarily sees
Gaussian statistics, because projection from high to low dimensions produces Gaussians.
This is a geometric statement: the shape of the wavefunction is fixed by the shape of
the arena.

\section{Orthogonal Measurement Bases}

\subsection{Slicing induces measurement}

In the cascade, each step selects an orthogonal direction and integrates over it. From the
perspective of a 4D observer, this operation has a natural interpretation: choosing a
direction along which to slice is choosing an observable to measure. The act of slicing---
projecting the $(d+1)$-ball onto the equatorial $d$-ball---discards information about the
perpendicular coordinate. This is measurement: extracting a definite value along one axis
at the cost of information about the complementary axis.

The compactification results of Section~\ref{sec:compactification} sharpen this
interpretation. Each slicing step does not annihilate the integrated-out direction; it
suppresses it to an effective radius $R_{\rm eff}(d)=1/\sqrt{d+3}$
(Theorem~\ref{thm:reff}). The weight function $(1-x^2)^{d/2}$ remains positive on $(-1,1)$
for all finite $d$. Measurement in the cascade is asymptotic suppression, not sharp
projection: the complementary coordinate is exponentially suppressed but never exactly zero.
This is a geometric realisation of the fact that quantum measurement does not destroy the
unmeasured degree of freedom---it renders it inaccessible at scale $R_{\rm eff}$.

\begin{theorem}[Orthogonal bases from slicing]\label{thm:bases}
The cascade's slicing structure, restricted to a 4D subspace, determines a family of
orthogonal bases for $\mathbb{R}^4$. Each choice of slicing axis within the 4D subspace
defines a measurement basis. Two measurements are complementary if and only if their slicing
axes are orthogonal in $\mathbb{R}^4$.
\end{theorem}
\begin{proof}
The slicing integral decomposes $B^{d+1}$ along a chosen axis $e$. Restricted to a 4D
subspace, the choice of $e$ selects a direction in $\mathbb{R}^4$. Different choices of $e$
give different decompositions of $\mathbb{R}^4$ into (axis)$\times$(3-ball). Two axes $e_1$,
$e_2$ are orthogonal if and only if the corresponding decompositions share no common axis,
which is the definition of complementary observables.
\end{proof}

\subsection{Non-commutativity from sequential slicing}

Slicing along axis $e_1$ then $e_2$ is not the same as slicing along $e_2$ then $e_1$,
because the first slice suppresses information about the first axis before the second slice
extracts information about the second. The cascade's slicing along axis $e$ extracts the
component $z = x\cdot e$ and suppresses the remainder. The extraction is an orthogonal
projection: $P_e(x) = (x\cdot e)\,e$, a rank-1 operator with $P_e^2 = P_e$.

\begin{theorem}[Commutator from slicing angle]\label{thm:commutator}
Let $e_1$, $e_2$ be two slicing axes in the 4D subspace with $e_1\cdot e_2 = \cos\theta$.
The commutator of the corresponding projections satisfies
\begin{equation}\label{eq:commutator}
\|P_{e_1}P_{e_2} - P_{e_2}P_{e_1}\| = \sin\theta\,|\!\cos\theta|.
\end{equation}
This vanishes if and only if $\theta = 0$, $\pi/2$, or $\pi$. The order-dependence of
sequential slicing is maximal at $\theta = \pi/4$.
\end{theorem}
\begin{proof}
The rank-1 projection along $e$ is $P_e = e\,e^T$, so
$(P_{e_1}P_{e_2})(x) = (x\cdot e_2)(e_1\cdot e_2)\,e_1$ and
$(P_{e_2}P_{e_1})(x) = (x\cdot e_1)(e_1\cdot e_2)\,e_2$. Writing $c = \cos\theta$:
\[
(P_{e_1}P_{e_2} - P_{e_2}P_{e_1})(x) = c\bigl[(x\cdot e_2)\,e_1 - (x\cdot e_1)\,e_2\bigr].
\]
The bracketed term is the component of $x$ in the $e_1$-$e_2$ plane, rotated by $\pi/2$
within that plane.  Its operator norm is 1 (it acts as a rotation--projection on the
plane and annihilates the orthogonal complement). Therefore
$\|P_{e_1}P_{e_2} - P_{e_2}P_{e_1}\| = |c|\cdot 1 = |\!\cos\theta|$... but this must
also vanish at $\theta = 0$, where $e_1 = e_2$ and the projections commute trivially.

More carefully: the bracketed operator $B(x) = (x\cdot e_2)\,e_1 - (x\cdot e_1)\,e_2$
restricted to the $e_1$-$e_2$ plane has matrix
$\bigl(\begin{smallmatrix}0&1\\-1&0\end{smallmatrix}\bigr)$
in the $(e_1,e_2)$ basis when $e_1\cdot e_2 = 0$, but this basis is not orthonormal
when $\theta\neq\pi/2$. In an orthonormal basis $\{e_1,\,e_\perp\}$ where
$e_2 = c\,e_1 + s\,e_\perp$ with $s = \sin\theta$:
\[
B(x) = (x\cdot e_2)\,e_1 - (x\cdot e_1)\,e_2
= s\bigl[(x\cdot e_\perp)\,e_1 - (x\cdot e_1)\,e_\perp\bigr].
\]
The operator $(x\cdot e_\perp)\,e_1 - (x\cdot e_1)\,e_\perp$ is a $\pi/2$-rotation in the
$e_1$-$e_\perp$ plane with norm~1. Therefore $\|B\| = |s| = |\!\sin\theta|$ and
\[
\|P_{e_1}P_{e_2} - P_{e_2}P_{e_1}\| = |c|\,|s| = |\!\sin\theta\cos\theta|
= \tfrac{1}{2}|\!\sin 2\theta|.
\]
This vanishes iff $\sin 2\theta = 0$, i.e., $\theta\in\{0,\pi/2,\pi\}$.
Maximum at $\theta = \pi/4$ (where $\sin 2\theta = 1$).
\end{proof}

\begin{remark}[Two distinct sources of commutativity]
The commutator vanishes at $\theta = 0$ and $\theta = \pi$ (parallel axes: same
measurement, trivially commuting) and at $\theta = \pi/2$ (orthogonal axes: the projections
have orthogonal ranges, so $P_{e_1}P_{e_2} = P_{e_2}P_{e_1} = 0$). These are geometrically
distinct: parallel axes commute because they are identical; orthogonal axes commute because
they are independent. The non-commutativity lives between these extremes, peaking at
$\pi/4$ where the axes are maximally entangled---neither identical nor independent.
\end{remark}

\begin{remark}[Complementarity vs non-commutativity]
In standard quantum mechanics, ``complementary observables'' (position and momentum) are
maximally non-commuting. In the cascade, orthogonal slicing axes ($\theta = \pi/2$) give
\emph{commuting} rank-1 projections. The non-commutativity of position and momentum arises
not from the angle between two single slicing axes, but from the relationship between a
\emph{complete basis} (all four axes simultaneously) and a rotated basis. For a pair of
orthonormal bases $\{e_j\}$ and $\{f_k\}$ related by angle $\theta$, the relevant
non-commutativity is between the complete projection operators
$\Pi_j = \sum_j \lambda_j P_{e_j}$ and $\Pi_k = \sum_k \mu_k P_{f_k}$, which do not
commute when the bases are related by a nontrivial rotation. The single-axis result
$\|[P_{e_1},P_{e_2}]\| = \tfrac{1}{2}|\!\sin 2\theta|$ is the elementary building block;
the full commutation relation $[x,p] = i\hbar$ requires the complexification of
Section~\ref{sec:complex} and the identification of $\hbar$ in Section~\ref{sec:hbar}.
\end{remark}

\section{The Born Rule from the Geometry of the Sphere}

\subsection{The question}

The cascade observer lives on $S^{d-1}$ (Section~2). Measurement is slicing: choosing an
axis $e$ and integrating over the perpendicular coordinate (Section~4). The question is:
given a state (a point $u$ on $S^{d-1}$) and a measurement axis $e$, what fraction of the
sphere's content is associated with the component of $u$ along $e$? The answer must come
from the geometry of the sphere alone. No probability postulate is assumed.

\subsection{Geometric content of a spherical cap}

\begin{theorem}[Spherical cap fraction]\label{thm:cap}
Let $u\in S^{d-1}$ be a unit vector and $e\in S^{d-1}$ a measurement axis with
$\cos\theta = u\cdot e$. The fraction of the uniform measure on $S^{d-1}$ contained in
the spherical cap $\{x\in S^{d-1} : x\cdot e \geq \cos\theta\}$ is
\begin{equation}\label{eq:cap}
F_d(\theta) = \frac{\int_0^{\theta}\sin^{d-2}\!\phi\,d\phi}
{\int_0^{\pi}\sin^{d-2}\!\phi\,d\phi}
= \frac{B(\sin^2\!\theta;\,\tfrac{d-1}{2},\,\tfrac{1}{2})}{B(\tfrac{d-1}{2},\,\tfrac{1}{2})},
\end{equation}
where $B(x;\,a,b)$ is the incomplete Beta function.
\end{theorem}
\begin{proof}
The uniform measure on $S^{d-1}$ in spherical coordinates with polar angle $\phi$ measured
from $e$ is $d\mu = \sin^{d-2}\!\phi\,d\phi\,d\Omega_{d-2}$. The angular part
$d\Omega_{d-2}$ integrates to $\Omega_{d-2}$ and cancels in the fraction. The remaining
integral is the stated Beta function ratio, via the substitution $t = \sin^2\!\phi$.
\end{proof}

\subsection{Concentration at \texorpdfstring{$\cos^2\theta$}{cos\^{}2 theta}}

\begin{theorem}[Born rule from concentration of measure]\label{thm:born}
As $d\to\infty$, the cap fraction $F_d(\theta)$ concentrates: for any state $u$ with
$u\cdot e = \cos\theta$, the probability that a uniformly random point $x\in S^{d-1}$
has $|x\cdot e|^2$ within $\varepsilon$ of $\cos^2\theta$ approaches~1 exponentially
in $d$. The unique consistent probability assignment for the outcome associated with
axis $e$ is
\[
p(e\mid u) = (u\cdot e)^2 = \cos^2\theta.
\]
This is the Born rule.
\end{theorem}
\begin{proof}
\textbf{Step 1} (marginal distribution). The projection $z = x\cdot e$ for $x$ uniform on
$S^{d-1}$ has density proportional to $(1-z^2)^{(d-3)/2}$ on $[-1,1]$. This is the
cascade's own slicing integrand at dimension $d-1$
(Lemma~\ref{lem:gaussian}): a Gaussian of width $\sigma = 1/\sqrt{d-1}$.

\textbf{Step 2} (concentration). By Lemma~\ref{lem:gaussian}, the density
$(1-z^2)^{(d-3)/2} \approx \exp(-(d-1)z^2/2)$ concentrates at $z = 0$ with width
$1/\sqrt{d-1}$. For a state at angle $\theta$ from $e$, decompose $x = z\,e + \sqrt{1-z^2}\,w$
where $w\in S^{d-2}$. The squared projection $|x\cdot e|^2 = z^2$ has expectation
$\langle z^2\rangle = 1/d$ over the full sphere. But conditioned on the state $u$ being at
angle $\theta$ from $e$, the relevant quantity is the conditional distribution of $z^2$
in the neighbourhood of $u$.

\textbf{Step 3} (the geometric argument). Consider the great circle through $u$ and $e$.
On this circle, $u$ has coordinate $z = \cos\theta$. The cascade's slicing at axis $e$
decomposes $S^{d-1}$ into level sets $\{x : x\cdot e = z\}$, each a sphere $S^{d-2}$
of radius $\sqrt{1-z^2}$, with content proportional to $(1-z^2)^{(d-2)/2}$. The content
at level $z = \cos\theta$ relative to the total is:
\begin{equation}\label{eq:born-ratio}
\frac{(1-\cos^2\theta)^{(d-2)/2}\,\delta z}
{\int_{-1}^{1}(1-z^2)^{(d-2)/2}\,dz}
= \frac{(\sin^2\theta)^{(d-2)/2}}{B(\tfrac{1}{2},\,\tfrac{d}{2})}.
\end{equation}
This peaks at $\theta = \pi/2$ (equator, maximum cross-section) and vanishes at $\theta = 0$
(pole). But the observer does not measure which level set $u$ lies on; the observer measures
the \emph{squared projection} $z^2 = \cos^2\theta$---the fraction of the unit vector's
length along the measurement axis. This quantity is determined by the angle alone.

\textbf{Step 4} (frame-function uniqueness via Cauchy additivity).
Let $P(e\mid u)$ denote the probability of the measurement outcome
associated with axis $e$ when the state is $u\in S^{d-1}$. By the
cascade's $SO(d)$-invariance of the uniform measure, $P(e\mid u)$
depends only on the inner product $u\cdot e$:
$P(e\mid u) = f(u\cdot e)$ for some continuous $f:[-1,1]\to[0,1]$
(continuity follows from the smoothness of the measure density
$(1-z^2)^{(d-3)/2}$). Invariance under $u\mapsto -u$---antipodal states
represent the same ray in the projective sphere $S^{d-1}/\mathbb{Z}_2$, and
the Pythagorean decomposition~$1 = \cos^2\theta + \sin^2\theta
= (u\cdot e)^2 + \|u-(u\cdot e)e\|^2$ is invariant under this
involution---forces $f(z) = g(z^2)$ for some continuous
$g:[0,1]\to[0,1]$.

For any orthonormal frame $\{e_1,\ldots,e_d\}$, the squared projections
$x_j = (u\cdot e_j)^2$ form a probability vector on the simplex $\Delta^{d-1}$:
\[
\sum_{j=1}^{d}(u\cdot e_j)^2 \;=\; \|u\|^2 \;=\; 1
\]
by Parseval's identity on $S^{d-1}$. Since a complete orthonormal frame
exhausts the measurement outcomes, the probabilities sum to~$1$:
\begin{equation}\label{eq:frame-function}
\sum_{j=1}^{d} g(x_j) \;=\; 1, \qquad (x_1,\ldots,x_d)\in\Delta^{d-1}.
\end{equation}
Every $(x_1,\ldots,x_d)\in\Delta^{d-1}$ is realised by some $(u,\{e_j\})$:
take $u = (\sqrt{x_1},\ldots,\sqrt{x_d})$ in the standard basis. Hence
\eqref{eq:frame-function} holds for \emph{every} probability vector in
$\Delta^{d-1}$.

Assume $d\geq 3$ (satisfied by the cascade's observer dimension~$d=4$;
Gleason's theorem fails for $d=2$, and the cascade's argument below is the
obstacle). For arbitrary $x,y\geq 0$ with $x+y\leq 1$, apply
\eqref{eq:frame-function} to $(x,y,1-x-y,0,\ldots,0)$ and to
$(x+y,1-x-y,0,\ldots,0)$ and subtract:
\[
g(x) + g(y) = g(x+y) + g(0).
\]
Setting $y=0$ forces $g(0) = 0$, so $g$ is Cauchy-additive on $[0,1]$.
Continuity of $g$ promotes additivity to linearity: $g(x) = \lambda x$ for
some $\lambda\geq 0$. Normalisation $g(1) = 1$ (certainty when $u = e_j$)
fixes $\lambda = 1$, so $g(x) = x$ uniquely.

Therefore $P(e\mid u) = g((u\cdot e)^2) = (u\cdot e)^2 = \cos^2\theta$,
with no other continuous function of $\theta$ satisfying
\eqref{eq:frame-function} on $\Delta^{d-1}$.
\end{proof}

\begin{remark}[Real-state special case of the complex Born rule]\label{rem:born-real-special-case}
The result $p(e\mid u) = (u\cdot e)^2$ uses the real inner product on
$S^{d-1}$. Under the $U(1)$ gauge forced by
Theorem~\ref{thm:U1-gauge-forced} (\S7.5), the cascade's actual Born
rule is the complex one,
\[
P(e\mid u) \;=\; |\langle e, u\rangle_{\mathbb{C}}|^2
\;=\; (u\cdot e)^2 + (Ju\cdot e)^2,
\]
which uniquely fixes the gauge-invariant quadratic functional under
$U(1)$. Theorem~\ref{thm:born} as stated derives the
\emph{real-state special case}: when $u$ and $e$ both lie in a real
subspace such that $J u$ is orthogonal to $e$ (equivalently, when
the two-dimensional subspace $\mathrm{span}_{\mathbb{R}}\{u, e\}$ is
$J$-invariant only trivially), the cross-term $(Ju\cdot e)^2$ vanishes
and the complex Born rule reduces to $(u\cdot e)^2$. The cascade's
slicing axes, taken at the cascade's natural per-layer pairing,
furnish such states. The derivation here is the cascade-internal
proof of the cosine-squared rule on the real-state slice; the
extension to general complex states is supplied by the gauge
structure of \S7.5 and Remark~\ref{rem:U1-gauge-Born}.
\end{remark}

\begin{remark}[What the proof uses]
The derivation uses three ingredients: (i)~the uniform measure on $S^{d-1}$ (from boundary
dominance, Section~2.3); (ii)~concentration of the slicing integrand $(1-z^2)^{(d-2)/2}$
(Lemma~\ref{lem:gaussian}); (iii)~Parseval's identity for the unit sphere
($\sum_j(u\cdot e_j)^2 = 1$). All three are geometric facts about the unit sphere. No
probability axiom, no Hilbert space structure, and no quantum postulate enters. The Born
rule is a theorem about spheres, not an axiom about measurement.
\end{remark}

\begin{remark}[Step 4 is the cascade-native frame-function uniqueness]\label{rem:sp15-status}
Step~4 derives the specific form $P(e\mid u) = \cos^2\theta$ from the
constraint $\sum_j g(x_j) = 1$ on $\Delta^{d-1}$ via Cauchy additivity
plus continuity plus normalisation. This is logically the
\emph{frame-function uniqueness} content of Gleason's
theorem~\cite{gleason}, specialised to real probability vectors on the
classical simplex rather than quantum projectors on a Hilbert
lattice. Three observations clarify the relationship:
\begin{enumerate}
\item \emph{What is derived cascade-natively.} The arena---Hilbert space
structure (Section~2), uniform measure on $S^{d-1}$ (boundary dominance),
concentration (Section~3), orthogonal decomposition and Parseval
(Section~4)---is constructed from the cascade's geometry before the
Cauchy argument runs. Gleason's antecedents are the cascade's theorems.
\item \emph{What the Cauchy argument adds.} Given the arena, the step
from the frame-function constraint \eqref{eq:frame-function} to
$g(x) = x$ uniquely is elementary: triple-point additivity on
$\Delta^{d-1}$ gives Cauchy's equation, continuity gives linearity,
normalisation fixes the slope. It does not require Gleason's full
machinery (measures on lattices of subspaces, complex projectors,
sub-additivity on incompatible observables); the cascade's state space
is $S^{d-1}/\mathbb{Z}_2$ and the Cauchy argument suffices.
\item \emph{The $d\geq 3$ condition.} The Cauchy argument requires a
third coordinate to apply the triple-point subtraction; $d=2$ admits
non-$z^2$ solutions ($g(x) + g(1-x) = 1$ alone has many continuous
solutions), which is exactly why Gleason's theorem fails at Hilbert
dimension~$2$. The cascade's observer sits at $d=4$ (Part~I), so
$d \geq 3$ is satisfied structurally, not assumed.
\end{enumerate}
The earlier phrasing ``concentration of measure forces this partition
to be the unique assignment consistent with additivity across
orthogonal axes'' elided the Cauchy step as intuition; the explicit
derivation above makes the logical structure visible.
\end{remark}

\begin{remark}[$U(1)$ gauge equivalence closes the
$\mathbb{R}P^3$ vs $\mathbb{C}P^1$ gap]\label{rem:U1-gauge-Born}
Step~4 derives the Born rule on the real simplex with the state space
identified as $S^{d-1}/\mathbb{Z}_2 = \mathbb{R}P^{d-1}$ (antipodal
quotient). Standard QM identifies the projective state space as
$\mathbb{C}P^{n-1}$ via $U(1)$ phase quotient, which differs from
the cascade's $\mathbb{R}P^{d-1}$ by the Hopf fibre dimension.

This gap closes once the cascade's complex structure $J$
(Theorem~\ref{thm:complex}) is recognised as a gauge equivalence
rather than merely a dynamical symmetry. Direct check: the cascade
discrete propagator $L(d) = i\,N(d)$ commutes with the $U(1)$ action
$R_\phi := \cos\phi\,I + \sin\phi\,J$ generated by $J$:
\[
[L(d),\,R_\phi] = 0 \quad\text{for all $\phi$.}
\]
Treating this $U(1)$ as a gauge equivalence (verified in
\texttt{tools/\allowbreak verifiers/\allowbreak cascade\_\allowbreak path\_\allowbreak state\_\allowbreak born\_\allowbreak rule.py})
reduces the projective state space from $\mathbb{R}P^{d-1}$ to
$\mathbb{C}P^{n-1}$ via the Hopf bundle $S^{2n-1} \to \mathbb{C}P^{n-1}$
($n = d/2$ for even~$d$).

Under the $U(1)$-gauge interpretation, Step~4's Cauchy argument runs
on the $U(1)$-quotient of the complex simplex:
$f(c) = g(|c|^2)$ by gauge invariance, and $\sum_i g(|c_i|^2) = 1$
on $\sum_i |c_i|^2 = 1$ forces $g(x) = x$ uniquely for $n \geq 3$
by the same triple-point Cauchy subtraction. The result is the
QM Born rule $P(e_i \mid \psi) = |\langle e_i, \psi\rangle_{\mathbb{C}}|^2$
identically. For $n = 2$ (qubit), single-system Cauchy is
under-determined exactly as in Gleason's theorem; closure follows
from compositionality with higher-dim cascade content (the
matter-content path tensor product of Part~IVa Section~4 reaches
$n \geq 3$ for any composite system through gauge layers).

The earlier $\mathbb{R}P^{d-1}$ identification corresponds to the
weaker gauge $\{1, -1\} = \mathbb{Z}_2 \subset U(1)$ (only antipodes
identified). Upgrading the gauge to the full $U(1)$
generated by $J$ recovers standard QM. Both choices give the
same numerical Born rule on real bases; the $U(1)$ choice
additionally recovers the QM complex Born rule
$|\langle u, v\rangle|^2$ as the
unique gauge-invariant quadratic functional, via the identity
$|\langle u, v\rangle_{\mathbb{C}}|^2 = (u\cdot v)^2 + (Ju\cdot v)^2$
(real inner product plus its $J$-rotated component).

\medskip
\noindent\textit{Resolution.} The above presents the $U(1)$ upgrade as
one of two cascade-mathematically-admissible gauge choices.
Theorem~\ref{thm:U1-gauge-forced} (\S7.5) closes the question:
$U(1)$ is forced under the cascade hypothesis, because $\mathbb{Z}_2$
gauge would predict an observable 2-cycle Planck-period oscillation
in Born-rule probabilities (from the cascade's scalar effective
Hamiltonian), absent from observation. See
Remark~\ref{rem:U1-closure-status} for the closure detail and the
empirical commitment it implies (the Renou falsifier).
\end{remark}

\subsection{What the cascade adds to Gleason}

Gleason's theorem~\cite{gleason} proves that the Born rule is the unique probability measure
on the lattice of subspaces of a Hilbert space of dimension $\geq 3$, but it assumes the
Hilbert space framework. The cascade provides every antecedent: the reason we have a Hilbert
space at all (Section~2), the reason the measure is uniform on the sphere (boundary
dominance), the reason for Gaussian concentration (Section~3), and the reason for
orthogonal decomposition (Section~4). Gleason proves uniqueness given the arena; the
cascade derives the arena, and Step~4's Cauchy functional-equation argument
(Remark~\ref{rem:sp15-status}) supplies the uniqueness step directly on the
real probability simplex without invoking Gleason's full frame-function
theorem.

\section{Complexification: The Forced Precession}
\label{sec:complex}

The cascade as described in~\cite{paperI} is real: it operates on $\mathbb{R}^d$. Quantum
mechanics requires complex amplitudes. This section shows how complex structure emerges from
the cascade, and---crucially---why the precession angle is not a free parameter.

\subsection{The forced precession}

\begin{theorem}[Forced precession]\label{thm:precession}
The angle between consecutive slicing axes is $\alpha = \pi/2$, forced by the cascade's
orthogonality axiom. No free parameter enters.
\end{theorem}
\begin{proof}
The cascade's orthogonality axiom is: each new slicing direction is perpendicular to all
previously integrated directions.

Applied to the slicing integral at step $d$: the slicing axis $e_{d+1}$ is perpendicular to
the equatorial hyperplane, forcing the half-integer argument in
$B(\frac{1}{2},\cdot)$, which gives $\Gamma(\frac{1}{2})=\sqrt{\pi}$ in the recurrence
(Theorem~3.1 of~\cite{paperI}).

Applied to consecutive axes: $e_{k+1}\perp e_k$ for all $k$ (the new axis is perpendicular
to the direction just integrated out), giving
$\alpha=\angle(e_k,e_{k+1})=\pi/2$.

Both conclusions follow from the
single axiom ``new slicing directions are perpendicular to all previously integrated
directions,'' applied at steps $d+1$ and $d$ respectively. The cascade's axiom forces both
simultaneously; there is no independent second assumption.

More precisely: the axiom is $e_k\perp e_j$ for all $k<j$ (each new slicing direction is
perpendicular to all earlier ones, which have been integrated out). For consecutive pairs
$(k,k+1)$: $\alpha=\angle(e_k,e_{k+1})=\pi/2$.
\end{proof}

\begin{corollary}[No free parameter in the precessing cascade]\label{cor:nofree}
The precessing cascade has no free geometric parameters. The precession angle $\alpha=\pi/2$
is determined by the cascade's orthogonality axiom alone.
\end{corollary}

\begin{remark}[Relationship to the $\sqrt{\pi}$ derivation]
Theorem~\ref{thm:precession} does not derive $\alpha=\pi/2$ from $\sqrt{\pi}$ or vice
versa. Both are consequences of the orthogonality axiom. The relationship is: the same
axiom ($e_k\perp$ all previously integrated directions) forces $\Gamma(\tfrac{1}{2})
=\sqrt{\pi}$ when applied to the integral over a single slicing direction, and forces
$\alpha=\pi/2$ when applied to the angle between consecutive slicing directions. They are
two theorems with a common hypothesis, not a chain where one implies the other.
\end{remark}

\subsection{Two quarter-turns yield the imaginary unit}

\begin{theorem}[Complex structure from forced precession]\label{thm:complex}
Let the cascade precess with $\alpha=\pi/2$ (forced by Theorem~\ref{thm:precession}). Then
the composition of two consecutive slicing operations, restricted to the plane of precession,
acts as multiplication by $-1$. The cascade's state space is naturally identified with a
complex Hilbert space, where the precession plane defines the complex structure $J$
($J^2=-\mathrm{Id}$).
\end{theorem}
\begin{proof}
A quarter-turn about $e_1$ maps $e_2\mapsto -e_1$ in the $e_1$-$e_2$ plane. A quarter-turn
about $e_2$ then maps $e_1\mapsto e_2$. The composition is a half-turn, i.e., multiplication
by $-1$. Define $J:e_1\mapsto e_2$, $e_2\mapsto -e_1$. Then $J^2=-\mathrm{Id}$, the
algebraic definition of a complex structure on $\mathbb{R}^2$. Extending to all pairs of
precessing axes partitions $\mathbb{R}^{2n}$ into $n$ complex dimensions $\mathbb{C}^n$,
with the cascade's quarter-turn providing the geometric origin of $i=e^{i\pi/2}$.

Since $\alpha=\pi/2$ is forced (Theorem~\ref{thm:precession}), this complex structure is
not introduced as an assumption but derived.
\end{proof}

\begin{remark}[Static angle to dynamical evolution: no hidden assumption]\label{rem:sp16-status}
Theorem~\ref{thm:precession} establishes a \emph{static} geometric fact
---the angle between consecutive slicing axes is $\pi/2$. The step to
``$J$ is the cascade's quantum evolution operator'' is often read as
promoting this static relation to a dynamical one, which would
introduce an assumption not contained in the static theorem. No such
assumption enters; the promotion is a composition of three
already-derived cascade ingredients:
\begin{enumerate}
\item \emph{Time is the slicing axis.} Section~7 below identifies the
observer's time coordinate with the cascade's slicing direction, on
the grounds that the slicing operation is irreversible and the
direction of irreversibility is the arrow of time. Under this
identification, the cascade's step from layer $d$ to layer $d-1$ is
one unit of temporal evolution.
\item \emph{Each evolution step is a quarter-turn.} Under~(1),
the operator $L(d)$ that advances the state by one cascade step
\emph{is} the quantum evolution operator over one time step.
Theorem~\ref{thm:precession} says each such step rotates consecutive
axes by $\pi/2$, so $L(d)$ carries the factor $e^{i\pi/2}=i$
(Section~7.3, Theorem~\ref{thm:discrete-schrodinger}).
\item \emph{Axis rotation lifts to state rotation via the passive
transformation law.} A rotation of the reference frame by $\theta$
acts on expansion coefficients of any state in that frame as the
inverse rotation, $-\theta$. The cascade's forced precession
rotates slicing axes by $\pi/2$ per step; the induced rotation on
states in the slicing basis is therefore $-\pi/2$, which is exactly
the complex-structure action $J^{-1}=-J$ on the $e_1$--$e_2$
plane. Lifting axis rotations to state rotations is not an
additional modelling step but the standard active/passive
transformation identity; no new cascade axiom enters.
\end{enumerate}
Hence ``$J$ evolves states'' is not a dynamical assumption but the
composition of (1) time-equals-slicing + (2) one-step quarter-turn +
(3) passive-transformation lift. Each ingredient is proved elsewhere
in the paper; the promotion is deductive. The earlier phrasing
``this complex structure is not introduced as an assumption but
derived'' is correct, and the derivation chain is now made
explicit.
\end{remark}

\begin{remark}[$U(1)$ generated by $J$ as gauge equivalence]\label{rem:U1-gauge-J}
The $J$ of Theorem~\ref{thm:complex} generates a continuous
$U(1) \subset SO(2n)$ acting on the cascade's $\mathbb{R}^{2n}$ state
space:
\[
R_\phi := \cos\phi\,I + \sin\phi\,J, \qquad \phi \in [0, 2\pi).
\]
This $U(1)$ is a dynamical \emph{symmetry} of the cascade evolution
($[L(d), R_\phi] = 0$ at every step, since
$L(d) = i\,N(d) = J\cdot N(d)$ commutes with itself and its powers).
Promoting it to a \emph{gauge equivalence}
(Remark~\ref{rem:U1-gauge-Born}) refines the projective state space
from $\mathbb{R}P^{d-1}$ to $\mathbb{C}P^{n-1}$ via the Hopf bundle
$S^{2n-1} \to \mathbb{C}P^{n-1}$, and the cascade Born rule under
gauge invariance becomes the QM complex Born rule
$|\langle u, v\rangle_{\mathbb{C}}|^2$. Within Part~II's per-step
content alone, $\mathbb{Z}_2$ and $U(1)$ are both
cascade-mathematically admissible.
Theorem~\ref{thm:U1-gauge-forced} (\S7.5) closes the question:
the cascade-internal cumulative phase $i^{N(t)-4}$ in the propagator
$K(N(t),4)$ would be observable under $\mathbb{Z}_2$ as a
Planck-period 2-cycle ray rotation (from a scalar effective
Hamiltonian), which conflicts with the cascade hypothesis; under
$U(1)$ the cumulative phase is gauge-trivial. The $U(1)$ choice is
therefore forced under hypothesis, not merely a structural
preference. See Remark~\ref{rem:U1-closure-status} for the closure
detail.
\end{remark}

\begin{remark}[The propagator phase and sphere topology at each cascade level]\label{rem:bott-phases}
The forced precession of $\pi/2$ per step accumulates a total propagator phase of
$(d-4)\cdot\pi/2$ from the observer at $d=4$ to level $d$. The phase has period~4. Each
layer $d$ also operates on the sphere $S^{d-1}$, whose topology depends on parity: $S^{d-1}$
is even-dimensional when $d$ is odd, and the hairy ball theorem forces every continuous
tangent field on an even-dimensional sphere to have a zero.
\begin{center}
\begin{tabular}{cccl}
\toprule
$d\bmod 4$ & Phase & $\chi(S^{d-1})$ & Cascade character \\
\midrule
0 & $+1$ & 0 & Real phase, no obstruction \\
1 & $i$  & 2 & Imaginary phase, forced zero \\
2 & $-1$ & 0 & Real phase, no obstruction \\
3 & $-i$ & 2 & Imaginary phase, forced zero \\
\bottomrule
\end{tabular}
\end{center}
Both columns are derived from the cascade's geometry. The phase column follows from
Theorem~\ref{thm:complex}: it is the accumulated precession $(d-4)\cdot\pi/2$ reduced
mod $2\pi$, with no external input. The sphere-topology column is elementary: $S^{d-1}$
is even-dimensional when $d$ is odd; even-dimensional spheres have Euler characteristic
$\chi = 2$, and the hairy ball theorem forces every tangent field to vanish; odd-dimensional
spheres have $\chi = 0$ and admit nonvanishing fields.
\end{remark}

\begin{corollary}[Phase--obstruction lockstep]\label{cor:lockstep}
The cascade's propagator phase is imaginary ($\pm i$) if and only if the sphere $S^{d-1}$
carries a forced tangent-field zero.
\end{corollary}
\begin{proof}
Both conditions reduce to the same parity test. The phase
$e^{i(d-4)\pi/2}$ is imaginary when $(d-4)$ is odd, i.e., when $d$ is odd. The sphere
$S^{d-1}$ is even-dimensional when $d$ is odd, and the hairy ball theorem forces a zero
precisely on even-dimensional spheres. Therefore: $d$ odd $\Leftrightarrow$ imaginary phase
$\Leftrightarrow$ forced zero. The forced precession and the hairy ball theorem are
geometrically independent---one concerns the angle between consecutive slicing axes, the
other concerns tangent fields on a sphere---but the orthogonality axiom forces them into
lockstep: the same parity of $d$ that gives an imaginary propagator phase also gives a
topological obstruction.
\end{proof}

\begin{remark}[Period 4 vs period 8]\label{rem:period4-vs-8}
The cascade's phase classification has period~4. The Clifford algebra
$\mathrm{Cl}(1,d-1)$ provides a finer classification with period~8 (Bott periodicity),
distinguishing layers that carry irreducibly complex spinor representations from those
that are purely real. This refinement---and its consequences for gauge structure and
fermion generations---is established in Part~III using the Clifford algebra classification
that also forces $d = 4$ via Lovelock uniqueness. At the level of this paper, the period-4
classification is complete: it derives the propagator's complex character and the
topological obstruction structure from the cascade's orthogonality axiom alone, without
importing the Clifford algebra or any framework from quantum field theory.
\end{remark}

\subsection{Interference}

Complex structure immediately gives interference. If two cascade paths arrive at the same
projected state with different accumulated precession angles $\varphi_1$ and $\varphi_2$,
their contributions are $e^{i\varphi_1}$ and $e^{i\varphi_2}$. The 4D observer measures:
\[
|e^{i\varphi_1}+e^{i\varphi_2}|^2 = 2+2\cos(\varphi_1-\varphi_2),
\]
exhibiting constructive interference when $\varphi_1=\varphi_2$ and destructive when
$\varphi_1-\varphi_2=\pi$.

\section{Time as Orthogonal Descent}

\subsection{The identification}\label{subsec:time-id}

We identify each cascade slicing operation as one elementary unit of time:
slicing \emph{is} dimensional collapse, and dimensional collapse \emph{is}
time. Each step of the cascade---integrating out one orthogonal direction
at a layer---resolves one more dimension of the unresolved infinite-
dimensional starting object $B^\infty$ into a definite cascade structure.
The cascade does not happen in time; the cascade's dimensional resolution
\emph{is} time.

\medskip
\noindent\textbf{Static slicing and dynamic collapse are the same
operation.} Read statically, the slicing recurrence
($V_{d+1} = \int V_d\,dx$ at $d+1$) is a structural decomposition of the
unit ball. Read dynamically, each application of the recurrence is one
Planck tick: one perpendicular direction is collapsed, the cascade's
truncation height advances by one layer, and the observer experiences one
elementary moment of time. The slicing operation is the dynamics; there is
no separate evolution operator acting on a pre-existing static structure.
Part~VI develops the cosmological consequences: each tick
adds one layer at the top of the truncated tower,
$N(t + \alpha\,t_{\mathrm{Pl,red}}) = N(t) + 1$, with the observer's
dimension fixed at $d = 4$ and the cascade's height growing above.

\medskip
\noindent\textbf{Geometric reading: asymptotic infall toward the
pseudo-horizon.} The cover-sheet thought experiment puts the observer on
the $S^3$ shell of a 5D black hole, asymptotically falling toward the
pseudo-horizon. Under the present identification, each slicing step
\emph{is} one asymptotic step of that infall: one dimension is collapsed
at the cascade's boundary, and the observer's $S^3$-host shell advances
one Planck tick closer to the pseudo-horizon. The infall never completes
because $B^\infty$ has infinitely many dimensions to resolve; the
asymptotic character of the cover-sheet thought experiment is exactly the
asymptotic character of the cascade's resolution of $B^\infty$. The
within-moment projection (Theorem~\ref{thm:propagator} below, $K(N(t), 4)
= \prod_{j=4}^{N(t)-1} L(j)$) is the cumulative effect of all collapses
performed up to time $t$, threaded back through the cascade to the
observer.

\medskip
\noindent\textbf{Three vocabularies, one rate.} The arrow of time, the
cascade's slicing direction, and the observer's asymptotic infall are
three readings of one fact: the rate at which the cascade resolves
$B^\infty$. Part~VI's tower-growth picture and the cover
sheet's pseudo-horizon picture are dual descriptions of this rate
(observer falls deeper $\leftrightarrow$ tower grows above), and Paper~II's
slicing direction is the cascade-internal ordering by which dimensions
are resolved. The ``forced direction'' of the thought experiment is the
forced direction of dimensional collapse.

\medskip
\noindent\textbf{Irreversibility.} The arrow of time is the direction of
dimensional collapse: the slicing integral loses information about the
integrated-out direction, and there is no inverse without supplying
external data the cascade does not contain. Time runs forward because
$B^\infty$ has not yet been fully resolved; backward time would require
``unresolving'' a dimension, which the cascade has no operation to
perform.

\subsection{The lapse function}

The cascade's lapse function is the ratio of adjacent volumes:
\[
N(d) = V_d/V_{d-1} = \sqrt{\pi}\cdot R(d),\qquad R(d) = \frac{\Gamma((d+1)/2)}{\Gamma((d+2)/2)}.
\]
Asymptotically $N(d)\approx\sqrt{2\pi/d}$ in
the Stirling regime. Properties: $N(d)\to 0$ as $d\to\infty$ (time does not exist at the
cascade's starting point); $N(d)$ monotonically increasing as $d$ decreases; $N(4)=3\pi/8
\approx 1.178$ (the 4D observer lives in a moderate regime; the Stirling approximation
$\sqrt{\pi/2}\approx 1.25$ is 6\% high).

\subsection{The discrete propagator}

\begin{theorem}[Cascade propagator]\label{thm:propagator}
The composition of cascade steps from $d=D$ down to $d=d'$ defines a propagator
$K(D,d')=\prod_{j=d'}^{D-1}L(j)$. In the precessing cascade, each $L(j)$ acquires the
forced phase factor $e^{i\alpha}=e^{i\pi/2}=i$, giving
$K(D,d')=\prod_j|L(j)|\cdot i^{D-d'}$.
\end{theorem}

\begin{remark}[Wilson-holonomy reading of $K$]\label{rem:wilson-K}
Equivalently, $K(D, d')$ is the (Abelian) Wilson holonomy of a
cascade-internal connection along the descent path from $d'$ to
$D-1$. The connection 1-form is the cascade slicing potential
$A(d) := p(d) = \tfrac{1}{2}\psi((d+1)/2) - \tfrac{1}{2}\ln\pi$
(introduced in Part~IVa as the cascade potential; here recognised
as the layer index 1-form on the descent path), and
\[
\bigl|K(D, d')\bigr|^2 \;=\; \exp\!\Bigl(2\sum_{j=d'}^{D-1} p(j)\Bigr),
\]
the gauge-invariant content of the holonomy. Under the $U(1)$
gauge interpretation
(Remark~\ref{rem:U1-gauge-J}), the modulus-squared is the physical
content; the phase $i^{D-d'}$ is the gauge-fixed representative of
the $U(1)$ Wilson-line phase. This re-reading does not change any
prediction; it identifies the cascade running coupling formulae of
Part~IVa Section~4 as Wilson holonomies in the multiplicative
scaling group, with $A(d) = p(d)$ derived cascade-internally from
the digamma function reflecting log sphere-area decay
($p(d) = -\tfrac{d}{dd}\log\Omega_d$). Verifier
\texttt{tools/verifiers/cascade\_wilson\_line.py}.
\end{remark}

\subsection{The Schr\"{o}dinger equation as effective description}

The cascade is discrete: 213 integer steps from $d=217$ to $d=4$. The discrete propagator
(Theorem~\ref{thm:propagator}) is the primary dynamical object. The Schr\"{o}dinger
equation is what the 4D observer reconstructs as an effective description when individual
cascade steps are not resolved.

\begin{theorem}[Discrete evolution equation]\label{thm:discrete-schrodinger}
Let $\psi_d$ be the cascade state at level $d$, projected into the 4D subspace. The
one-step evolution is
\[
\psi_{d-1} = L(d)\,\psi_d, \qquad L(d) = N(d)\cdot e^{i\pi/2} = i\,N(d),
\]
where $N(d) = \sqrt{\pi}\cdot R(d)$ is the lapse (Section~7.2) and the phase $i$ is forced
by Theorem~\ref{thm:complex}. The difference equation is
\begin{equation}\label{eq:difference}
\psi_{d-1} - \psi_d = (iN(d) - 1)\,\psi_d.
\end{equation}
\end{theorem}
\begin{proof}
Direct substitution of $L(d) = iN(d)$ into the one-step relation.
\end{proof}

\begin{lemma}[Lapse normalisation]\label{lem:lapse-norm}
Given the discrete evolution $\psi_{d-1} = iN(d)\,\psi_d$ of
Theorem~\ref{thm:discrete-schrodinger}, define the cumulative-lapse-normalised
wavefunction at descent reference $D \geq d$ by
\[
\tilde\psi_d \;:=\; \psi_d \cdot \prod_{j=d+1}^{D} N(j)^{-1}.
\]
Then $\tilde\psi$ evolves by a pure phase per cascade step:
\[
\tilde\psi_{d-1} \;=\; i\,\tilde\psi_d,
\qquad
|\tilde\psi_{d-1}| = |\tilde\psi_d|.
\]
The $|L(j)| = N(j)$ factor in the discrete propagator
(Theorem~\ref{thm:propagator}) is fully absorbed into the normalisation
constants $\prod_j N(j)^{-1}$; what remains is the precessing phase
$i^{D-d}$, unitary by construction.
\end{lemma}
\begin{proof}
Direct substitution:
\[
\tilde\psi_{d-1}
= \psi_{d-1} \prod_{j=d}^{D} N(j)^{-1}
= iN(d)\,\psi_d \cdot N(d)^{-1} \prod_{j=d+1}^{D} N(j)^{-1}
= i\,\tilde\psi_d.
\]
The modulus identity $|\tilde\psi_{d-1}| = |i||\tilde\psi_d| = |\tilde\psi_d|$
follows.
\end{proof}

\begin{corollary}[Effective Schr\"{o}dinger equation]\label{cor:schrodinger}
Let the 4D observer parametrise the cascade descent by a coordinate $t$ with
$\Delta t = N(d)$ per step (proper time from the lapse). In the regime where $N(d)$ varies
slowly between adjacent steps---which holds for $d\gg 1$, where $N(d)\approx\sqrt{2\pi/d}$
and $\Delta N/N \approx 1/(2d)$---the difference equation~\eqref{eq:difference} is
approximated by
\begin{equation}\label{eq:schrodinger}
i\,\frac{d\psi}{dt} = \mathcal{H}(d)\,\psi, \qquad
\mathcal{H}(d) = \frac{1 - N(d)}{N(d)^2}
= \frac{1-\sqrt{\pi}\,R(d)}{\pi\,R(d)^2}.
\end{equation}
\end{corollary}
\begin{proof}
Apply the lapse normalisation of Lemma~\ref{lem:lapse-norm}: pass from
$\psi_d$ to $\tilde\psi_d = \psi_d\prod_{j>d}^D N(j)^{-1}$, so that the
modulus-changing cascade factor is fully stored in the normalisation
constants and $\tilde\psi$ evolves by a pure phase per step,
$\tilde\psi_{d-1} = i\tilde\psi_d$. Dividing the difference equation for
$\tilde\psi$ by $\Delta t = N(d)$:
\[
\frac{\tilde\psi_{d-1} - \tilde\psi_d}{N(d)}
= \frac{i - 1}{N(d)}\,\tilde\psi_d.
\]
Multiplying both sides by $-i$ and rearranging:
\[
i\,\frac{\tilde\psi_d - \tilde\psi_{d-1}}{N(d)}
= \frac{1 - N(d)}{N(d)^2}\cdot N(d)\,\tilde\psi_d
\;+\; \text{terms subleading in the slowly-varying regime.}
\]
In the $\Delta N/N = O(1/d)$ regime, replacing the difference quotient
by $d\tilde\psi/dt$ yields the Schr\"{o}dinger
form~\eqref{eq:schrodinger} with
$\mathcal{H}(d) = (1-N(d))/N(d)^2$ on the normalised wavefunction
$\tilde\psi$. Dropping the tilde (the normalisation is a choice of
reference and does not affect the Schr\"{o}dinger operator) gives the
stated form.
\end{proof}

\begin{remark}[Derivation status of the Schr\"{o}dinger reduction]\label{rem:sp18-status}
The effective Schr\"{o}dinger equation~\eqref{eq:schrodinger} factorises
into three ingredients with three different derivation statuses:
\begin{enumerate}
\item \emph{Derived.} The discrete evolution
$\psi_{d-1} = iN(d)\psi_d$ (Theorem~\ref{thm:discrete-schrodinger})
and the exact discrete propagator
$K(D,d') = \prod_j|L(j)|\cdot i^{D-d'}$
(Theorem~\ref{thm:propagator}). These are exact cascade facts, not
approximations.
\item \emph{Derived.} The unitary reduction via lapse normalisation
(Lemma~\ref{lem:lapse-norm}): the cumulative-lapse product
$\prod_{j>d}N(j)^{-1}$ is a specific construction, not an informal
absorption; $\tilde\psi$ evolves by a unit-modulus phase $i$ per step,
so any Schr\"{o}dinger operator extracted from the $\tilde\psi$
evolution is automatically Hermitian in the continuum limit.
\item \emph{Approximate (effective description).} The continuum limit
replacing the difference quotient by $d\tilde\psi/dt$, and the
specific functional form $\mathcal{H} = (1-N)/N^2$, are effective at
$d\gg 1$ where $\Delta N/N = O(1/d)$. At the observer's
$d = 4$, $N(4) = 3\pi/8 = 1.178$ and the fractional change per
step is $\sim 12\%$, so the continuum form is a coarse-grained
reconstruction. The exact cascade dynamics is the discrete propagator
(Theorem~\ref{thm:propagator}); the differential equation is the
observer's effective description.
\end{enumerate}
The earlier phrasing ``the imaginary part contributes a
decay/growth that is absorbed into the lapse normalisation'' named
this construction without exhibiting it; Lemma~\ref{lem:lapse-norm}
makes the absorption explicit as the cumulative product
$\prod_{j>d} N(j)^{-1}$, which is a definite function of $d$ and $D$
rather than a conventional choice.
\end{remark}

\begin{remark}[Status of the approximation]
The continuum limit is valid deep in the cascade ($d\gg 1$), where consecutive lapses
differ by $O(1/d)$. At the observer's dimension $d=4$, $N(4) = 3\pi/8 = 1.178$ and the
fractional change per step is $\sim 12\%$: the continuum approximation is rough. The
4D observer's effective Schr\"{o}dinger equation is therefore a coarse-grained description
of a discrete dynamics, not an exact law. This is consistent with the cascade's structure:
the discrete propagator (Theorem~\ref{thm:propagator}) is exact; the differential equation
is the observer's reconstruction of it.
\end{remark}

\begin{remark}[The Hamiltonian in cascade units]
$\mathcal{H}(d) = (1-N(d))/N(d)^2$ is a Gamma function ratio: it is $\mathcal{H}(d) =
(1-\sqrt{\pi}\,R(d))/(\pi\,R(d)^2)$, where $R(d) = \Gamma((d+1)/2)/\Gamma((d+2)/2)$.
At $d = 4$: $\mathcal{H}(4) = (1-3\pi/8)/(3\pi/8)^2 = -0.128$. At $d=19$:
$\mathcal{H}(19) = 2.10$. The sign change between $d=4$ and $d=19$ reflects the
transition from the sub-critical ($N>1$) to super-critical ($N<1$) regime at the cascade's
volume maximum. The value of $\hbar$ enters when converting $\mathcal{H}$ to physical
energy units (Section~\ref{sec:hbar}).
\end{remark}

\begin{remark}
The arrow of time is structural, not thermodynamic. The cascade descends because the slicing
integral compresses $(d+1)\to d$; no inverse operation restores the integrated-out direction
without supplying external information.
\end{remark}

\subsection{The \texorpdfstring{$U(1)$}{U(1)} gauge is forced}

\begin{theorem}[$U(1)$ gauge forced by scalar-Hamiltonian consistency]\label{thm:U1-gauge-forced}
Under the cascade hypothesis (Section~1: the infinite-dimensional unit
ball, descended to four dimensions, is indistinguishable from our universe),
the cascade's gauge group on the projective state space is the full $U(1)$
generated by $J$, not the proper subgroup
$\mathbb{Z}_2 = \langle J^2\rangle = \{\pm I\}$.
\end{theorem}

\begin{proof}
Two candidate gauge groups are cascade-internally distinguishable. From
Theorem~\ref{thm:complex}, $J^2 = -\mathrm{Id}$, so $\langle J^2\rangle =
\{I, -I\} = \mathbb{Z}_2$ is a cascade-derived subgroup; from the same
theorem, $\langle J\rangle$ generates the full $U(1) \subset SO(2n)$ via
$R_\phi := \cos\phi\,I + \sin\phi\,J$. Both commute with the cascade's
per-step propagator (Remark~\ref{rem:U1-gauge-J}: $[L(d), R_\phi] = 0$ for
all $\phi$). The question is which is the correct gauge equivalence on
the projective state space.

By Theorem~\ref{thm:discrete-schrodinger}, the per-step propagator is
\[
L(d) \;=\; i\,N(d) \;=\; J \cdot N(d),
\]
which acts as a complex scalar on every state vector: a pure phase $i = J$
multiplied by a real positive scalar $N(d)$, with no operator content
distinguishing different state-space directions. By
Corollary~\ref{cor:schrodinger}, the corresponding effective Schr\"odinger
equation has scalar Hamiltonian $\mathcal{H}(d) = (1-N(d))/N(d)^2$, again
with no operator structure on the state space's dimensions.

Apply $L(d)$ to a generic state $\psi$ and pass to projective space under
each candidate gauge:
\begin{itemize}
\item Under $\mathbb{Z}_2$: rays are equivalence classes
$[\psi] = \{\psi, -\psi\}$ in $\mathbb{R}P^{d-1}$. The map
$\psi \mapsto i\psi = J\psi$ sends $[\psi]$ to $[J\psi]$, a different
ray in $\mathbb{R}P^{d-1}$ for generic $\psi$ (since $J\psi \neq \pm\psi$).
The per-step propagator therefore moves rays in a 2-cycle:
$[\psi]\to [J\psi]\to[\psi]\to\ldots$, with successive Born-rule
probabilities $\langle\psi, e\rangle_R^2$ and $\langle J\psi, e\rangle_R^2$
at any measurement axis $e$. The cascade's per-step propagator generates
observable per-step ray rotation from a scalar Hamiltonian.
\item Under $U(1)$: rays are equivalence classes
$[\psi] = \{e^{i\phi}\psi : \phi \in [0, 2\pi)\}$ in $\mathbb{C}P^{n-1}$.
The map $\psi \mapsto i\psi$ sends $[\psi]$ to $[i\psi] = [\psi]$ (same
ray, since $i \in U(1)$ is gauged out). The per-step propagator acts as
the identity on rays; the cumulative cascade phase $i^{D-d'}$ in
$K(D, d') = \prod|L(j)| \cdot i^{D-d'}$ (Theorem~\ref{thm:propagator})
is gauge-trivial, and only the real factor $\prod|L(j)|$ has projective
content.
\end{itemize}

Standard quantum mechanics has the property that a scalar Hamiltonian
generates only an unobservable global phase on rays:
$e^{-i\mathcal{H}t}\psi$ and $\psi$ are the same physical state. The
cascade's effective Schr\"odinger equation
(Corollary~\ref{cor:schrodinger}) is mathematically identical in form to
standard Schr\"odinger evolution; under the cascade hypothesis the
cascade's prediction must therefore agree with standard QM on the
physical content of scalar-Hamiltonian evolution. The $\mathbb{Z}_2$
gauge predicts observable per-step ray rotation from the scalar
$\mathcal{H}(d)$, contradicting standard QM and therefore contradicting
the cascade hypothesis. The $U(1)$ gauge predicts no per-step ray
rotation, matching standard QM. Therefore the cascade's gauge group is
$U(1)$.
\end{proof}

\begin{remark}[Closure of the gauge-upgrade question]\label{rem:U1-closure-status}
Theorem~\ref{thm:U1-gauge-forced} closes the question of whether the
$U(1)$ upgrade of Remarks~\ref{rem:U1-gauge-Born} and~\ref{rem:U1-gauge-J}
is forced or merely permitted: it is forced, conditional on the cascade
hypothesis, at the same level as every other Tier 1--2 result of the
cascade series (each of which derives observation-matching predictions
from the hypothesis without further empirical input). The closure is
internal to Part~II: it uses only Theorem~\ref{thm:complex} (existence
of $J$ and $J^2$), Theorem~\ref{thm:discrete-schrodinger} (scalar
per-step propagator $L(d) = iN(d)$), Theorem~\ref{thm:propagator} (the
cumulative cascade propagator), and Corollary~\ref{cor:schrodinger}
(scalar effective Hamiltonian).

The earlier framing in Remark~\ref{rem:U1-gauge-Born} (``both gauge
choices are cascade-consistent; the $U(1)$ choice additionally matches
standard QM'') and Remark~\ref{rem:U1-gauge-J} (``the choice is
structural'') is now superseded: $\mathbb{Z}_2$ is cascade-mathematically
admissible at the per-step level, but its predicted observable 2-cycle
ray rotation under scalar $\mathcal{H}(d)$ is excluded by the cascade
hypothesis. Numerical verifier:
\texttt{tools/verifiers/u1\_gauge\_closure\_verifier.py} (V5).

\medskip
\noindent\textbf{Empirical commitment: the Renou falsifier.}
Theorem~\ref{thm:U1-gauge-forced} commits the cascade to
complex-amplitude QM (with complex Born rule
$|\langle u, v\rangle_\mathbb{C}|^2$) as opposed to real-amplitude QM
(with real Born rule $\langle u, v\rangle_\mathbb{R}^2$). Renou et
al.~\cite{renou2021} (\emph{Nature} 600, 625--629, 2021,
``Quantum theory based on real numbers can be experimentally
falsified'') showed that real-amplitude QM and complex-amplitude QM
predict measurably different statistics in tripartite Bell scenarios.
A future tripartite Bell experiment yielding real-amplitude statistics
would therefore falsify the cascade as derived in Part~II. Conversely,
the experimental evidence to date (Renou \emph{et al.} demonstrate
violation of the real-QM bound in a photonic implementation) is
consistent with the cascade's complex-amplitude commitment.
\end{remark}

\subsection{Cascade dynamics: kinematic closure}\label{subsec:dynamics-closure}

The identification ``one cascade slicing operation $=$ one Planck tick
$=$ one dimensional collapse'' is now established at three independent
levels of the series:
\begin{itemize}[nosep]
\item \emph{Cover sheet} (``One Infinity, Not Many''). Time is the rate
at which the resolution of $B^\infty$ proceeds; each Planck tick
resolves one more dimension of $B^\infty$ into a definite cascade
layer.
\item \emph{Prelude}, ``Resolution as cosmic time'' remark, ``three
vocabularies for one rate'': mathematical (slicing recurrence),
physical (Planck tick), geometric (asymptotic infall toward the
pseudo-horizon).
\item \emph{Section~\ref{subsec:time-id} above}: ``slicing $=$
dimensional collapse $=$ time''; static and dynamic readings of the
slicing operation are the same operation viewed through different
vocabularies.
\end{itemize}
This subsection records that, under this shared identification, the
cascade's \emph{kinematic} dynamics is fully specified by the discrete
propagator (Theorem~\ref{thm:propagator}) plus the tower-growth
operator of Paper~VI~\cite{paper6}. The ``missing Hamiltonian''
objection --- that the cascade lacks a state-dependent generator of
time evolution --- misreads where the cascade's dynamical structure
sits. The \emph{interaction} dynamics (gauge-coupled fermion couplings
with non-trivial state-space action) is structurally separate; its
cascade-native action form is determined by six conditions in
Part~IVb~\cite{paper4b}, Remark on the cascade-native gauge-coupled
fermion action
($\mathit{rem{:}fermion\text{-}gauge\text{-}coupling\text{-}proposal}$).

\medskip\noindent\textbf{The cascade has a c-number Hamiltonian, by
construction.} Corollary~\ref{cor:schrodinger} extracts the effective
Hamiltonian
\[
\mathcal{H}_{\rm eff}(d) \;=\; \frac{1 - N(d)}{N(d)^2}
\;=\; \frac{1 - \sqrt{\pi}\,R(d)}{\pi\,R(d)^2},
\]
a Gamma-function ratio that is a \emph{c-number}, not a state-dependent
operator. This is the cascade's analogue of the QM Hamiltonian:
state-independent because the cascade's dynamics is generated by the
slicing recurrence itself --- a geometric per-tick rate --- not by an
operator on the state space. Standard QM postulates a state-dependent
$\hat H$ from which $|\psi(t)\rangle = e^{-i\hat H t}|\psi(0)\rangle$
follows; the cascade derives the state-independent rate
$\mathcal{H}_{\rm eff}(d)$ from which
$|\psi_{d-1}\rangle = i\,N(d)\,|\psi_d\rangle$ follows
(Theorem~\ref{thm:discrete-schrodinger}). Both produce unitary
evolution; the cascade's version is per-tick and geometric, while QM's
is per-instant and state-space. The c-number character is not a defect
to be repaired by introducing an operator-valued generator; it is the
cascade's structural commitment that single-particle dynamics is
intrinsic to the geometry, not chosen via an action principle.

\medskip\noindent\textbf{The period-1 tick reading is forced
cascade-internally.} Three candidate identifications relate one
Planck tick to one cascade slicing step: period~1 (one tick $=$ one
$d$-collapse), period~2 (one tick $=$ one $\mathbb{Z}_2$ chirality
flip), period~4 (one tick $=$ one $J$-cycle returning to identity).
Two are excluded:
\begin{itemize}[nosep]
\item \emph{Period~2 excluded by gauge equivalence.}
Theorem~\ref{thm:U1-gauge-forced} forces the projective gauge to be
$U(1)$ generated by $J$, not the proper subgroup
$\mathbb{Z}_2 = \langle J^2\rangle$. Under $\mathbb{Z}_2$ the per-tick
phase $i$ would generate observable 2-cycle ray rotation from a
scalar Hamiltonian, contradicting standard QM; identifying the
period-2 chirality flip with one Planck tick would reinstate exactly
this contradiction, so the period-2 reading is structurally
forbidden.
\item \emph{Period~4 excluded cosmologically.} Under the period-4
reading, Paper~VI's pre-Big-Bang e-fold count
(Proposition~3.3 of~\cite{paper6}) divides by 4, giving
$N_e = 70.94/4 \approx 17.7$, well below the observational inflation
window $[50, 70]$ and falsified at the inflation-window level.
\item \emph{Period~1 (one Planck tick $=$ one $d$-collapse) survives.}
The per-tick phase $i$ is gauge-trivial under $U(1)$; the geometric
modulus per tick is $|L(j)| = N(j)$; Paper~VI's $\alpha = 1$ tick
rate gives $N_e = 70.94$, inside the inflation window. The cascade's
period-1 reading is therefore not chosen but forced.
\end{itemize}

\medskip\noindent\textbf{Cosmic-time evolution from tower growth.}
Paper~VI~\cite{paper6} promotes the truncation height $N(t)$ to a
physical time coordinate ($N(t + t_{\rm Pl,red}) = N(t) + 1$ under the
forced $\alpha = 1$). The cascade state evolves between moments by
\[
|\Psi(t + t_{\rm Pl,red})\rangle
\;=\; K(N(t)+1, 4) \circ A_{N(t)+1} \circ K(N(t), 4)^{-1}\,
|\Psi(t)\rangle,
\]
where $K$ is the within-moment discrete propagator (this section) and
$A_{N(t)+1}$ is the tower-growth operator that adds layer $N(t)+1$ at
the top. Both ingredients are cascade-internal: $K$ is forced by the
slicing recurrence; $A$ is the realisation of ``one tick $=$ one
$d$-collapse'' at the cosmic scale. Single-particle within-moment
dynamics and cosmic-scale between-moments dynamics are operationally
the same evolution mechanism applied to different slices of the
cascade tower.

\medskip\noindent\textbf{What is closed and what is open.}
The kinematic dynamics is closed. The cascade's free-particle and
cosmic-scale evolution operators are explicit
($K \circ A \circ K^{-1}$), the per-tick phase is gauge-trivial under
the forced $U(1)$, the per-tick modulus is the geometric lapse
$N(d)$, the c-number Hamiltonian is $\mathcal{H}_{\rm eff}(d) = (1-N)/N^2$,
and the period-1 tick reading is forced. Nothing is missing from the
single-particle evolution at the observer's frame; the discrete
propagator plus tower-growth supplies the full kinematic structure.

The interaction dynamics is open in a known place. Standard Model
scattering amplitudes, decay rates, and QCD running across scales
require a \emph{state-dependent} interaction Hamiltonian --- the
cascade analogue of $\bar\psi\gamma^\mu A_\mu\psi$. The
gauge-coupled fermion action
(Part~IVb~\cite{paper4b}, \emph{rem:fermion-gauge-coupling-proposal}:
$S_f = \sum_d m(d)\bar\psi\psi + \sum_{d\in\{12,13,14\}} g(d)\bar\psi
T^a A^a\psi$ with $m(d) = g(d) = \sqrt{\alpha(d)}$) is structurally
determined by six cascade-internal conditions (F1)--(F6); per-layer
Berezin reduction at $A = 0$ closes; the originally-listed downstream
items --- integer charge numerators in the $Y$ spectrum, photon
self-energy screening in $1/\alpha_{\rm em}$, and $\mathrm{SU}(2)_L$
parity-violation structure --- are now also closed via
Part~IVb~\cite{paper4b} Theorem on sector-fundamental U(1)$_Y$
($\mathit{thm{:}sector\text{-}fundamental\text{-}y}$),
Open Question on $1/\alpha_{\rm em}$ screening
($\mathit{oq{:}alpha\text{-}em\text{-}screening}$, resolved), and
Part~IVa~\cite{paper4a} Section~2 Spin$(12)$ Dirac decomposition
respectively.

The dynamical objection ``the cascade has no Lagrangian and therefore
cannot be a complete physical theory'' applies to the
interaction-dynamics gap, but that gap is now also substantially
closed at the cascade-native level: the action form is determined,
the downstream observables are derived, and per CLAUDE.md Check~7
no semiclassical Green's-function residual is admissible. The
kinematic dynamics is closed by this section; the interaction
dynamics is structurally complete with all originally-listed
load-bearing consequences derived.

\section{Compactification and the Geometry of Temporal Descent}
\label{sec:compactification}

Section~7 identifies time with slicing and states that the arrow of time arises from the
irreversibility of integrating out a direction. This section gives that identification
concrete geometric content: each slicing step does not eliminate a direction but
asymptotically suppresses it, with an exact compactification radius determined by the Beta
function.

\subsection{The effective radius of each temporal step}

\begin{theorem}[Compactification radius]\label{thm:reff}
At each slicing step $d$, the integrated-out direction retains an effective radius
\[
R_{\rm eff}(d) = \frac{1}{\sqrt{d+3}}.
\]
\end{theorem}
\begin{proof}
The effective radius is the RMS extent of the integrated-out coordinate under the
slicing measure:
\[
R_{\rm eff}^2 = \langle x^2\rangle
= \frac{\int_{-1}^{1} x^2(1-x^2)^{d/2}\,dx}{\int_{-1}^{1}(1-x^2)^{d/2}\,dx}.
\]
Substituting $t = x^2$, the numerator is $B(3/2,\,d/2+1)$ and the denominator is
$B(1/2,\,d/2+1)$. Their ratio is
\[
\langle x^2\rangle = \frac{B(3/2,\,d/2+1)}{B(1/2,\,d/2+1)}
= \frac{\Gamma(3/2)}{\Gamma(1/2)}\cdot\frac{\Gamma(d/2+3/2)}{\Gamma(d/2+5/2)}
= \frac{1}{2}\cdot\frac{1}{d/2+3/2}
= \frac{1}{d+3},
\]
where the first factor is $\Gamma(3/2)/\Gamma(1/2) = 1/2$ and the second uses
$\Gamma(x+1) = x\,\Gamma(x)$ once.
Therefore $R_{\rm eff} = 1/\sqrt{d+3} \to 0$ as $d\to\infty$.
\end{proof}

The compactification is asymptotic: $R_{\rm eff}\neq 0$ at any finite $d$. The weight
function $(1-x^2)^{d/2}$ is smooth and positive on $(-1,1)$ for all finite $d$. There is
no sharp boundary, no horizon, no topology change---only exponential suppression. The word
``asymptotic'' is essential: the compactification limit is approached but never reached at
any finite step.

\begin{center}
\begin{tabular}{cccc}
\toprule
$d$ & $R_{\rm eff}$ & $N(d)$ & Physical role \\
\midrule
4   & 0.37796 & 1.17810 & Observer dimension \\
5   & 0.35355 & 1.06667 & $S^3$ boundary layer \\
7   & 0.31623 & 0.91429 & Area maximum $d_0$ \\
19  & 0.21320 & 0.56755 & First threshold $d_1$ \\
217 & 0.06742 & 0.16997 & Second threshold $d_2$ \\
\bottomrule
\end{tabular}
\end{center}

\subsection{The scale coincidence}

The cascade has two natural length scales at each step: the compactification radius
$R_{\rm eff}(d) = 1/\sqrt{d+3}$ from the Beta function, and the Gaussian width
$\sigma(d) = 1/\sqrt{d}$ of the slicing integrand (Lemma~\ref{lem:gaussian}). These
are defined by different objects---one by the second moment, one by the integrand's
curvature at $x = 0$---yet they coincide asymptotically.

\begin{theorem}[Scale coincidence]\label{thm:scale}
\[
\frac{R_{\rm eff}(d)}{\sigma(d)} = \sqrt{\frac{d}{d+3}} \to 1 \quad\text{as }
d\to\infty,
\]
with correction $O(1/d)$: $R_{\rm eff}/\sigma = 1 - \frac{3}{2d} + O(d^{-2})$.
\end{theorem}
\begin{proof}
Direct computation from Lemma~\ref{lem:gaussian} and Theorem~\ref{thm:reff}.
\end{proof}

\begin{corollary}[Second universal cascade constant]\label{cor:erf}
The fraction of the slicing integrand's content within the compactification radius
satisfies
\[
F(d) = \frac{\int_{|x|<R_{\rm eff}(d)}(1-x^2)^{d/2}\,dx}
{\int_{-1}^{1}(1-x^2)^{d/2}\,dx}
\to \operatorname{erf}\!\left(\frac{1}{\sqrt{2}}\right) = 0.68269\ldots
\quad\text{as }d\to\infty,
\]
a second universal constant of the cascade alongside $\sqrt{\pi}$, determined entirely
by the coincidence $R_{\rm eff}/\sigma\to 1$.
\end{corollary}
\begin{proof}
In the Gaussian limit (Lemma~\ref{lem:gaussian}), the integrand is
$e^{-dx^2/2}$ with width $\sigma = 1/\sqrt{d}$. The fraction within
$R_{\rm eff} = 1/\sqrt{d+3}$ is
$\operatorname{erf}(R_{\rm eff}/(\sigma\sqrt{2}))$. As $d\to\infty$,
$R_{\rm eff}/(\sigma\sqrt{2})\to 1/\sqrt{2}$, giving
$\operatorname{erf}(1/\sqrt{2})$.
\end{proof}

\begin{remark}[The cascade's effective geometric origin]
After $d_2 = 217$, sphere areas satisfy $\Omega_d < 10^{-120}$---geometrically
indistinguishable from zero. The cascade's second threshold $d_2 = 217$ and its formal
starting point $d = \infty$ are therefore the same geometric point: the cascade begins at
geometric nothing, and nothing remains after $d_2$. The limit $d\to\infty$ in
Corollary~\ref{cor:erf} is consequently exact, not merely asymptotic, when applied to the
cascade tower: there is no geometric content beyond $d_2 = 217$ to distinguish the two
points. The constant $\operatorname{erf}(1/\sqrt{2})$ is a property of the cascade's
origin, evaluated at its effective terminus.
\end{remark}

\begin{remark}[Structural foreshadowing]
The cascade produces exactly two universal constants from its geometry:
$\sqrt{\pi} = \Gamma(1/2)$ from orthogonality, and $\operatorname{erf}(1/\sqrt{2}) =
0.68269$ from the scale coincidence. The first generates the cosmological constant.
The second is within $0.15\%$ of $(\pi-1)/\pi = 0.68169$, which Part~V derives as the
dark energy fraction $\Omega_\Lambda$. Both constants are properties of the
Gamma function evaluated at its cascade-distinguished points; neither is fitted.
\end{remark}

\subsection{Irreversibility as asymptotic suppression}

The identification of time with slicing (Section~7.1) asserts that the slicing integral
loses information about the integrated-out direction. Theorem~\ref{thm:reff} makes this
precise: the direction is not lost but suppressed to scale $R_{\rm eff}=1/\sqrt{d+3}$.
From the perspective of the $(d+1)$-dimensional geometry, the direction's contribution is
exponentially localised near $x=0$ with Gaussian width $\sigma\approx 1/\sqrt{d}$
(Lemma~\ref{lem:gaussian}). From the perspective of the $d$-dimensional residual geometry,
the direction has been compactified.

The arrow of time acquires a quantitative meaning: each step of temporal descent reduces one
direction's effective extent by a factor that depends on $d$. Early in the cascade (large
$d$), $R_{\rm eff}$ is small and the suppression per step is mild. Late in the cascade
(small $d$), $R_{\rm eff}$ is larger and each step represents a more substantial geometric
change. At the observer's dimension $d=4$, $R_{\rm eff}=1/\sqrt{7}\approx 0.378$: each
temporal step compactifies a direction to roughly 38\% of the unit ball's radius.

The irreversibility of slicing is now geometrically transparent. To reverse a slicing step
would require decompactifying a direction from $R_{\rm eff}$ back to unit radius. This
requires external information: the specific correlations between the integrated-out
coordinate and the remaining geometry, which are asymptotically suppressed in the slicing
integral. The arrow of time is the asymptotic approach to the compactification limit,
viewed from the descending side.

\subsection{Unitarity and information}

The cascade dynamics are unitary: each step $L(d) = iN(d)$ has $|L(d)| = N(d) > 0$, and
the discrete propagator $K(D,d')$ (Theorem~\ref{thm:propagator}) preserves the phase
structure. Yet each slicing step appears irreversible (Section~7.1). These two statements
are compatible because they describe different levels.

The full cascade geometry from $d=D$ down to $d=d'$ is described by a unitary propagator
$K(D,d')$ (Theorem~\ref{thm:propagator}). No information is lost in the total evolution.
A single slicing step, viewed from the $d$-dimensional residual geometry, appears to lose
information because the observer cannot access the compactified direction at scale
$R_{\rm eff}=1/\sqrt{d+3}$. The lost information is not destroyed; it is encoded in the
correlations between the compactified direction and the boundary.

This is made precise by boundary dominance: $\Omega_{d-1}/V_d=d$ (Theorem~3.1
of~\cite{paperI}). The boundary $S^{d-1}$ carries a fraction $d/(d+1)$ of the content at
each layer. Information that appears lost from the interior is retained on the boundary.
The cascade is unitary as a whole; it looks irreversible from any single layer because the
boundary of that layer encodes information the interior observer cannot directly access.

\begin{remark}[Comparison with the measurement problem]
The structure is suggestive: unitary global evolution, apparent irreversibility from a
restricted perspective, information retained on the boundary. These are the ingredients of
both the quantum measurement problem and the black hole information paradox. The cascade
provides them from geometry alone. Whether this structural parallel constitutes a resolution
of either problem is addressed in Section~\ref{sec:open}.
\end{remark}

\subsection{Geometric decoherence from compactification}

The cascade's slicing step at dimension $d$ integrates over the fibre coordinate $x$ with
weight $(1-x^2)^{d/2}$. This integration traces out the fibre degree of freedom, producing
a quantum channel on the base state space. A superposition of base states that are
distinguished by their fibre coordinates loses coherence through this tracing. The
decoherence rate is set by the compactification radius $R_{\rm eff}$.

\begin{theorem}[Geometric decoherence rate]\label{thm:decoherence}
Let $|\Psi\rangle = (|a\rangle + |b\rangle)/\sqrt{2}$ be a superposition of two states on
$S^d$, where $a$ and $b$ have fibre coordinates $x_a = a\cdot e_{d+1}$ and
$x_b = b\cdot e_{d+1}$ along the slicing axis. After tracing over the fibre with the
cascade's slicing weight, the reduced density matrix is
\[
\rho_{\rm red} = \tfrac{1}{2}\bigl(w_a|\bar{a}\rangle\langle\bar{a}|
+ w_b|\bar{b}\rangle\langle\bar{b}|
+ D\,|\bar{a}\rangle\langle\bar{b}|
+ D^*\!|\bar{b}\rangle\langle\bar{a}|\bigr),
\]
where $\bar{a},\bar{b}$ are the base projections, $w_a,w_b$ are slicing weights, and the
decoherence factor is
\begin{equation}\label{eq:decoherence}
D = \frac{(1-x_a x_b)^{d/2}}{\bigl[(1-x_a^2)(1-x_b^2)\bigr]^{d/4}}.
\end{equation}
For small fibre separation $\Delta x = |x_a - x_b|$ near the equator ($x_a,x_b\approx 0$):
\begin{equation}\label{eq:decoherence-gaussian}
D \approx \exp\!\left(-\frac{d\,(\Delta x)^2}{4}\right)
= \exp\!\left(-\frac{(\Delta x)^2}{4\,R_{\rm eff}^2}\right),
\end{equation}
with decoherence length $\ell_{\rm dec} = 2\,R_{\rm eff}(d) = 2/\sqrt{d+3}$.
\end{theorem}
\begin{proof}
The cascade's slicing decomposes $S^d$ into level sets $\{x\cdot e_{d+1}=z\}$, each a
sphere $S^{d-1}(\sqrt{1-z^2})$ with content proportional to $(1-z^2)^{(d-1)/2}$. A point
$a\in S^d$ at height $x_a$ has base projection $\bar{a}$ on $S^{d-1}(\sqrt{1-x_a^2})$.
The tracing operation replaces each state by its base projection, weighted by the slicing
amplitude $(1-x^2)^{d/4}$ at its height.

The off-diagonal element $\langle\bar{a}|\rho_{\rm red}|\bar{b}\rangle$ arises from the
overlap of the fibre amplitudes:
\[
D = \frac{\int_{-1}^{1}(1-z^2)^{d/2}\,\varphi_a^*(z)\,\varphi_b(z)\,dz}
{\sqrt{\int|\varphi_a|^2(1-z^2)^{d/2}dz\cdot\int|\varphi_b|^2(1-z^2)^{d/2}dz}},
\]
where $\varphi_a(z)$ and $\varphi_b(z)$ are the fibre wave functions. For states localised
at heights $x_a$ and $x_b$ with cascade-natural width $\sigma = 1/\sqrt{d}$
(Lemma~\ref{lem:gaussian}), the wave functions are
$\varphi_a(z)\propto\exp(-(z-x_a)^2 d/2)$ and
$\varphi_b(z)\propto\exp(-(z-x_b)^2 d/2)$.

The product $\varphi_a^*\varphi_b\propto\exp(-d[(z-x_a)^2+(z-x_b)^2]/2)
= \exp(-d(z-\bar{x})^2)\exp(-d(\Delta x)^2/4)$ where $\bar{x}=(x_a+x_b)/2$.
The factor $\exp(-d(\Delta x)^2/4)$ is independent of $z$ and survives the integration.
The remaining Gaussian integral over $z$, combined with the slicing weight, contributes a
ratio that approaches 1 for $x_a,x_b\approx 0$ (the Gaussian approximation of
Lemma~\ref{lem:gaussian} is centred at $z=0$). The result
is~\eqref{eq:decoherence-gaussian}.

The exact expression~\eqref{eq:decoherence} follows from the same calculation without the
Gaussian approximation. For $\Delta x\to 0$: $D\to 1$ (full coherence). For
$\Delta x \gg 2/\sqrt{d}$: $D\to 0$ (complete decoherence).
\end{proof}

\begin{corollary}[Decoherence length equals compactification radius]\label{cor:dec-length}
The cascade's decoherence length is $\ell_{\rm dec} = 2\,R_{\rm eff}$. States separated by
less than $R_{\rm eff}$ in the fibre direction retain coherence ($D > e^{-1/4}\approx 0.78$);
states separated by more than $2\,R_{\rm eff}$ are effectively decohered
($D < e^{-1}\approx 0.37$).
\end{corollary}

\begin{center}
\renewcommand{\arraystretch}{1.2}
\begin{tabular}{cccc}
\toprule
$d$ & $R_{\rm eff}$ & $\ell_{\rm dec}$ & $D$ at $\Delta x = R_{\rm eff}$ \\
\midrule
4   & 0.378 & 0.756 & 0.86 \\
7   & 0.316 & 0.632 & 0.88 \\
19  & 0.213 & 0.427 & 0.92 \\
217 & 0.067 & 0.135 & 0.97 \\
\bottomrule
\end{tabular}
\end{center}

\begin{remark}[Comparison with environment-induced decoherence]
Standard decoherence theory~\cite{zurek} gives a decoherence factor
$D = \exp(-\Delta x^2/\lambda_{\rm env}^2)$ where $\lambda_{\rm env}$ is the environment's
coherence length (the thermal de~Broglie wavelength for a thermal bath). The cascade
reproduces this structure exactly, with the identification:
\begin{center}
\begin{tabular}{ll}
\toprule
Standard decoherence & Cascade \\
\midrule
System & Base state on $S^{d-1}$ \\
Environment & Fibre coordinate $x$ \\
Tracing over environment & Slicing integral $\int(1-x^2)^{d/2}dx$ \\
Environmental coherence length $\lambda_{\rm env}$ & $2\,R_{\rm eff} = 2/\sqrt{d+3}$ \\
Decoherence factor $e^{-\Delta x^2/\lambda^2}$ & $e^{-(\Delta x)^2(d+3)/4}$ \\
\bottomrule
\end{tabular}
\end{center}
The cascade does not model the environment; it \emph{is} the environment. The
fibre coordinate that is integrated out at each slicing step is the geometric
degree of freedom that the lower-dimensional observer cannot access. The
Gaussian form of the decoherence factor is not assumed---it is derived from
the Gaussian concentration of the slicing integrand
(Lemma~\ref{lem:gaussian}). The decoherence length $2\,R_{\rm eff}$ is not
fitted---it is the compactification radius of Theorem~\ref{thm:reff},
determined by the Beta function.
\end{remark}

\subsection{Cumulative decoherence over the cascade descent}

Each of the 213 slicing steps traces out one fibre, contributing an overlap deficit
$1-C^2_{d,d+1}$ to the total decoherence. The cumulative structure is computable exactly
from the Beta function.

\begin{theorem}[Cumulative overlap deficit]\label{thm:cumulative-dec}
The total adjacent-layer overlap deficit over the full cascade descent is
\begin{equation}\label{eq:cum-deficit}
\sum_{d=5}^{216}\bigl(1-C^2_{d,d+1}\bigr) = 0.02108,
\end{equation}
where each $C_{d,d+1} = B(\tfrac{1}{2},d+\tfrac{3}{2})/
\sqrt{B(\tfrac{1}{2},d+1)\,B(\tfrac{1}{2},d+2)}$.
The sum converges: $72.7\%$ accumulates by $d_1=19$, and $96.8\%$ by $d=100$.
\end{theorem}
\begin{proof}
Each term is computed from the Beta function (Theorem~\ref{thm:decoherence}).
The asymptotic scaling $1-C^2\sim 1/(8d^2)$ gives convergence of
$\sum 1/(8d^2)<\infty$. Numerical evaluation at each integer $d$ gives the stated sum
(computed in \texttt{tools/model\_checks/cascade\_decoherence.py}).
\end{proof}

The full Gram correlation matrix for the 213-layer path $d=5,\ldots,217$ has eigenvalue
structure:
\begin{center}
\renewcommand{\arraystretch}{1.2}
\begin{tabular}{ccc}
\toprule
Eigenvalue & Value & Fraction of trace \\
\midrule
$\lambda_1/n$ & $0.9389$ & $93.89\%$ \\
$\lambda_2/n$ & $0.0541$ & $5.41\%$ \\
$\lambda_3/n$ & $0.0062$ & $0.62\%$ \\
$\lambda_4/n$ & $0.0008$ & $0.08\%$ \\
\midrule
$\varepsilon = 1-\lambda_1/n$ & \multicolumn{2}{c}{$0.0611 = 6.11\%$} \\
\bottomrule
\end{tabular}
\end{center}

\begin{corollary}[Cascade coherence]\label{cor:cascade-coherence}
The cascade retains $93.9\%$ coherence over the full 213-step descent. The remaining
$6.1\%$ is the inter-layer coupling: geometric content shared between layers rather than
propagating independently through each step. This eigenvalue deficit $\varepsilon$ is
the same quantity that drives the descent corrections in the Part~0 Supplement.
\end{corollary}

\begin{remark}[Two decoherence numbers]
The adjacent sum $\sum(1-C^2_{\rm adj}) = 0.021$ and the eigenvalue deficit
$\varepsilon = 0.061$ measure different things. The adjacent sum counts nearest-neighbour
losses only. The eigenvalue deficit counts all pairwise off-diagonal structure. Their ratio
$0.021/0.061 = 0.34$: long-range correlations between non-adjacent layers contribute twice
as much total decoherence as the adjacent steps alone.
\end{remark}

\begin{remark}[Separation-dependent decoherence]
For a superposition with angular separation $\Delta\theta$ distributed democratically
among all 217 dimensions ($\Delta x_j^2 = \Delta\theta^2/217$ per integrated-out
direction), the total decoherence factor is
\[
D(\Delta\theta) = \exp\!\Bigl(-\Delta\theta^2\sum_{d=5}^{217}\frac{d+3}{4\cdot 217}\Bigr)
= \exp(-28.0\,\Delta\theta^2).
\]
Coherence survives for $\Delta\theta\lesssim 0.19\,\mathrm{rad}\approx 11^\circ$ and is
destroyed for $\Delta\theta\gtrsim 30^\circ$.
\end{remark}

\section{The Role of \texorpdfstring{$\hbar$}{hbar}}
\label{sec:hbar}

\subsection{\texorpdfstring{$\hbar$}{hbar} as the unit-matching constant}

By Theorem~\ref{thm:precession}, the precession angle is fixed: $\alpha=\pi/2$. The cascade
is purely geometric and produces dimensionless quantities. Physical quantum mechanics has a
dimensionful constant $\hbar$ with units of action. It enters as:
\[
\hbar = \frac{\alpha}{\Delta t} = \frac{\pi/2}{\Delta t},
\]
where $\Delta t$ is the physical duration assigned to one cascade step. Since $\alpha$ is
now fixed by geometry, $\hbar$ encodes only the identification of cascade steps with
physical time increments. It is a unit-matching constant, not a free geometric parameter.

\subsection{The classical limit}

The classical limit $\hbar\to 0$ corresponds to $\Delta t\to\infty$ at fixed
$\alpha=\pi/2$: the observer's temporal resolution is too coarse to track the precession.
Phases decorrelate, and the 4D observer sees incoherent (classical) projections. This is
geometric decoherence.

\section{Tensor Products and Entanglement}

\subsection{Tensor structure from iterated slicing}

The slicing recurrence decomposes $V_{d+1}=V_d\times(\text{fibre})$, where the fibre is
the one-dimensional integral over the perpendicular coordinate. Iterating $k$ times:
$V_{d+k}=V_d\times\text{fibre}_1\times\cdots\times\text{fibre}_k$. This is a tensor
product: the $(d+k)$-dimensional space factors into a $d$-dimensional base and $k$
one-dimensional fibres.

\subsection{Entanglement}

\begin{theorem}[Generic entanglement]\label{thm:entangle}
For $d\geq 2$, the cascade state after one slicing step is generically entangled between
the base and fibre subsystems. Separable states have measure zero in the cascade's state
space.
\end{theorem}
\begin{proof}
The slicing cross-section at height $x$ has radius $\sqrt{1-x^2}$, creating a correlation
between the fibre coordinate $x$ and the base radius. A product state would require the
cross-sectional radius to be independent of $x$, which occurs only for $d=0$. For $d\geq 2$,
the coupling term $(1-x^2)^{d/2}$ is non-trivial, and the set of separable states is a
lower-dimensional submanifold.
\end{proof}

The compactification interpretation (Section~\ref{sec:compactification}) sharpens the
entanglement structure. The fibre coordinate $x$ is compactified to effective radius
$R_{\rm eff}=1/\sqrt{d+3}$, but the base radius $\sqrt{1-x^2}$ depends on $x$ throughout
the compactification region. The base--fibre entanglement is generated by this $x$-dependent
coupling and persists at all scales above $R_{\rm eff}$. Below $R_{\rm eff}$, the fibre is
effectively frozen and the entanglement is inaccessible to a $d$-dimensional observer---but
it is retained in the boundary correlations.

\subsection{Bell inequality violation}

The cascade's entangled states (Theorem~\ref{thm:entangle}) produce correlations that no
assignment of definite values can reproduce. This subsection proves it by explicit
computation on $S^7$.

\begin{lemma}[Classical CHSH bound]\label{lem:chsh}
For any four numbers $A,A',B,B'\in\{-1,+1\}$:
$|AB - AB' + A'B + A'B'|\leq 2$.
\end{lemma}
\begin{proof}
$AB - AB' + A'B + A'B' = A(B-B')+A'(B+B')$. Since $B,B'\in\{-1,+1\}$, exactly one of
$|B-B'|$ and $|B+B'|$ equals 2 and the other equals 0. Therefore $|S|=2$.
\end{proof}

\noindent Consequence: any model that assigns definite outcomes $\pm 1$ to all four
settings simultaneously---regardless of which pair is measured---has
$|\langle S\rangle|\leq 2$.

\begin{theorem}[CHSH violation on $S^7$]\label{thm:bell}
There exists a state $u\in S^7$ and measurement settings such that the CHSH correlator,
computed from the cascade's Born rule (Theorem~\ref{thm:born}), satisfies $S = 2\sqrt{2}$.
\end{theorem}
\begin{proof}
\textbf{Step 1} (bipartite structure). The complex structure $J$
(Theorem~\ref{thm:complex}) on $\mathbb{R}^8$ gives four complex coordinates
$\zeta_k = x_{2k-1}+ix_{2k}$ for $k=1,\ldots,4$. Partition into subsystem
$A=\{\zeta_1,\zeta_2\}$ and subsystem $B=\{\zeta_3,\zeta_4\}$. The tensor product
structure (Section~10.1) identifies $\mathbb{R}^8 = \mathbb{C}^4 =
\mathbb{C}^2_A\otimes\mathbb{C}^2_B$, with basis $\{e_1,e_3,e_5,e_7\}$ corresponding to
the real parts of $\{\zeta_1\zeta_3,\;\zeta_1\zeta_4,\;\zeta_2\zeta_3,\;\zeta_2\zeta_4\}$.

\textbf{Step 2} (entangled state). Take
$u = (e_1+e_7)/\sqrt{2}\in S^7$,
which has $\zeta_1=\zeta_4=1/\sqrt{2}$, all others zero. This state is not separable:
a product state $f(\zeta_1,\zeta_2)\otimes g(\zeta_3,\zeta_4)$ requires
$\zeta_1\zeta_4 = \zeta_2\zeta_3$, but $\tfrac{1}{2}\neq 0$.

\textbf{Step 3} (measurement directions). Alice measures at angle $\alpha$ via the axis
\[
a(\alpha) = \cos\tfrac{\alpha}{2}\,e_1 + \sin\tfrac{\alpha}{2}\,e_3\in\mathbb{R}^4_A
\]
(a unit vector in the $\zeta_1$-$\zeta_2$ plane). Bob measures at angle $\beta$ via
$b(\beta) = \cos\tfrac{\beta}{2}\,e_5 + \sin\tfrac{\beta}{2}\,e_7\in\mathbb{R}^4_B$.
The joint direction for outcome $(+,+)$ is the tensor product embedded in $\mathbb{R}^8$:
\[
v(\alpha,\beta) = \cos\tfrac{\alpha}{2}\cos\tfrac{\beta}{2}\,e_1
+ \cos\tfrac{\alpha}{2}\sin\tfrac{\beta}{2}\,e_3
+ \sin\tfrac{\alpha}{2}\cos\tfrac{\beta}{2}\,e_5
+ \sin\tfrac{\alpha}{2}\sin\tfrac{\beta}{2}\,e_7,
\]
a unit vector in $S^7$ ($|v|^2 = 1$).

\textbf{Step 4} (Born rule). By Theorem~\ref{thm:born} applied to axis $v$ on $S^7$:
\[
P(+,+) = (u\cdot v)^2 = \left[\tfrac{1}{\sqrt{2}}\cos\tfrac{\alpha}{2}\cos\tfrac{\beta}{2}
+ \tfrac{1}{\sqrt{2}}\sin\tfrac{\alpha}{2}\sin\tfrac{\beta}{2}\right]^2
= \tfrac{1}{2}\cos^2\!\tfrac{\alpha-\beta}{2}.
\]
Replacing $\alpha\to\alpha+\pi$ (Alice's $-$ outcome) and/or $\beta\to\beta+\pi$ (Bob's):
\[
P(+,-) = P(-,+) = \tfrac{1}{2}\sin^2\!\tfrac{\alpha-\beta}{2},\qquad
P(-,-) = \tfrac{1}{2}\cos^2\!\tfrac{\alpha-\beta}{2}.
\]
Check: $\sum P = 1$. The four outcome directions are orthogonal in $\mathbb{R}^8$ and
$u$ lies in their span, so the probabilities exhaust the Born rule by Parseval
(Theorem~\ref{thm:born}, Step~4).

\textbf{Step 5} (correlation). The correlation function is
\[
E(\alpha,\beta) = P(\text{agree})-P(\text{disagree})
= \cos^2\!\tfrac{\alpha-\beta}{2} - \sin^2\!\tfrac{\alpha-\beta}{2}
= \cos(\alpha-\beta).
\]

\textbf{Step 6} (CHSH). Choose $\alpha_1=0$, $\alpha_2=\pi/2$, $\beta_1=\pi/4$,
$\beta_2=3\pi/4$:
\begin{align*}
S &= E(0,\,\tfrac{\pi}{4}) - E(0,\,\tfrac{3\pi}{4})
+ E(\tfrac{\pi}{2},\,\tfrac{\pi}{4}) + E(\tfrac{\pi}{2},\,\tfrac{3\pi}{4})\\
&= \cos\bigl(-\tfrac{\pi}{4}\bigr) - \cos\bigl(-\tfrac{3\pi}{4}\bigr)
+ \cos\bigl(\tfrac{\pi}{4}\bigr) + \cos\bigl(-\tfrac{\pi}{4}\bigr)\\
&= \tfrac{1}{\sqrt{2}} - \bigl(-\tfrac{1}{\sqrt{2}}\bigr) + \tfrac{1}{\sqrt{2}}
+ \tfrac{1}{\sqrt{2}}
= \tfrac{4}{\sqrt{2}} = 2\sqrt{2}.\qedhere
\end{align*}
\end{proof}

\begin{remark}[Real-state special case under the forced gauge]\label{rem:bell-real-special-case}
The proof above uses the real Born rule $P(+,+) = (u\cdot v)^2$ with
state $u = (e_1+e_7)/\sqrt{2}$ and measurement directions $v$ in
$\mathrm{span}\{e_1, e_3, e_5, e_7\}$. Under the $U(1)$ gauge forced by
Theorem~\ref{thm:U1-gauge-forced} (\S7.5), the cascade's Born rule is
the complex one $|\langle u, v\rangle_{\mathbb{C}}|^2$. For the
particular states chosen in this proof, $Ju = (e_2+e_8)/\sqrt{2}$ has
no overlap with $\mathrm{span}\{e_1, e_3, e_5, e_7\}$, so
$Ju\cdot v = 0$ and the complex Born rule collapses to
$(u\cdot v)^2$ identically: the result $S = 2\sqrt{2}$ is independent
of the gauge choice for this particular state and these particular
measurement directions. The CHSH violation derived here is therefore
gauge-agnostic for the chosen real states; it is the special case in
which real and complex Born rules coincide. A general (complex)
cascade state would obey the complex Born rule
$|\langle u, v\rangle_{\mathbb{C}}|^2 = (u\cdot v)^2 + (Ju\cdot v)^2$
forced by \S7.5. Verifier:
\texttt{tools/verifiers/u1\_gauge\_closure\_verifier.py} (V3, V5).
\end{remark}

\begin{remark}[Bipartition of $\mathbb{C}^4$: cascade order fixes it; Tsirelson makes it moot]\label{rem:sp17-status}
Step~1 of the proof above takes the bipartition
$A = \{\zeta_1,\zeta_2\}$, $B = \{\zeta_3,\zeta_4\}$. Among the three
distinct 2+2 partitions of $\mathbb{C}^4$ --- namely
$(1,2)|(3,4)$, $(1,3)|(2,4)$, $(1,4)|(2,3)$ --- the cascade's
sequential slicing fixes which one is natural, and Tsirelson's bound
makes the choice numerically irrelevant.

\medskip\noindent\textbf{Cascade-order forcing of the partition.}
The slicing recurrence produces axes in a \emph{total order}
$e_1, e_2, e_3, \ldots$ by layer index: axis $e_k$ is introduced at
cascade step $k$ (Theorem~\ref{thm:precession}; the slicing-axis
$=$~time identification is proved in \S7 below). The complex structure
$J$ (Theorem~\ref{thm:complex}) pairs adjacent axes into complex
dimensions $\zeta_k = x_{2k-1} + ix_{2k}$, so the cascade order on
complex coordinates is $\zeta_1 < \zeta_2 < \zeta_3 < \zeta_4$ by
the order of their real-axis introduction. The only 2+2 bipartition
that respects this total order as ``first half vs.\ second half'' is
$A = \{\zeta_1,\zeta_2\}$, $B = \{\zeta_3,\zeta_4\}$. The
alternatives $(1,3)|(2,4)$ and $(1,4)|(2,3)$ interleave the order
and are not cascade-sequential.

\medskip\noindent\textbf{Tsirelson's bound is partition-independent.}
Tsirelson's theorem bounds the CHSH correlator for any maximally
entangled state on $\mathbb{C}^2\otimes\mathbb{C}^2$ at $2\sqrt{2}$,
regardless of which bipartition of a higher-dimensional space
produces the 2-qubit structure. For each of the three 2+2 partitions
of $\mathbb{C}^4$ above, the maximally entangled state in the chosen
bipartition saturates the bound at the same numerical value
$2\sqrt{2}$. The cascade's conclusion $S = 2\sqrt{2}$ is therefore
robust against the bipartition choice: cascade order forces the
natural partition to be $(1,2)|(3,4)$, but had the cascade happened
to suggest another, the Tsirelson value would still be $2\sqrt{2}$.
The bipartition is cascade-derived via order; the violation magnitude
is partition-independent.
\end{remark}

\begin{remark}[What the proof uses]
The ingredients are: (i)~a unit vector $u\in S^7$ (the cascade's state space at $d=8$);
(ii)~the complex structure $J$ from Theorem~\ref{thm:complex}, providing the tensor product
decomposition; (iii)~the Born rule $p=(u\cdot v)^2$ from Theorem~\ref{thm:born}, applied to
tensor product directions on $S^7$; (iv)~the classical CHSH bound from elementary
combinatorics. The violation $2\sqrt{2}>2$ is a theorem about unit vectors and inner products
on the 7-sphere.
\end{remark}

\begin{remark}[Why the cascade violates the classical bound]
The classical bound assumes that Alice's and Bob's outcomes are determined by a shared
assignment of definite values to all four settings simultaneously. The cascade state
$u=(e_1+e_7)/\sqrt{2}$ distributes its unit norm across two tensor product directions:
it has amplitude in both the $\zeta_1$ and $\zeta_4$ sectors, with no amplitude in
$\zeta_2$ or $\zeta_3$. No single pair of definite values for Alice and Bob can
reproduce this distribution across all measurement choices. The correlations arise from
the geometry of $S^7$, which the 4D observer cannot access directly---only through
projections that obey the Born rule.
\end{remark}

\section{Summary: From Cascade to Quantum Mechanics}

\begin{center}
\small
\begin{tabular}{>{\raggedright\arraybackslash}p{6cm}>{\raggedright\arraybackslash}p{6.5cm}c}
\toprule
QM Postulate & Cascade Origin & Section \\
\midrule
State space is a Hilbert space & Unit sphere from cascade geometry & 2 \\
States are unit vectors & Normalisation from ball boundary & 2 \\
Arena is the sphere, not the ball & Boundary dominance $\Omega/V=d$ & 2.3 \\
Gaussian wavefunctions & Projection from high dimensions (CLT) & 3 \\
Gaussian integrand $(1-x^2)^{d/2}$ & Concentration of slicing measure & 3.1 \\
Normalisation by $\sqrt{\pi}$ & Orthogonal compression constant & 3 \\
Orthogonal measurement bases & Slicing axis choices & 4 \\
Measurement as suppression & $R_{\rm eff}=1/\sqrt{d+3}$ & 4.1, 8 \\
Non-commutativity $\|[P,Q]\|=\tfrac{1}{2}|\sin 2\theta|$ & Rank-1 projection commutator & 4 \\
Born rule $p = \cos^2\theta$ & Sphere geometry + Parseval + concentration & 5 \\
Complex amplitudes & Forced precession $\alpha=\pi/2\Rightarrow J^2=-\mathrm{Id}$ & 6 \\
No free geometric parameters & $\alpha=\pi/2$ forced; same axiom as $\sqrt{\pi}$ & 6.1 \\
Phase--obstruction lockstep & Imaginary phase $\Leftrightarrow$ forced zero (parity of $d$) & 6.5 \\
Interference & Phase accumulation from forced precession & 6.3 \\
Time as dimensional collapse & One slicing $=$ one $d$-collapse $=$ one Planck tick & 7.1 \\
Unitary time evolution & Cascade propagator (discrete) & 7 \\
Schr\"{o}dinger equation & Effective description from discrete propagator & 7.4 \\
$U(1)$ gauge on $\mathbb{C}P^{n-1}$ forced; complex Born rule & Scalar $\mathcal{H}(d)$ + cascade hypothesis & 7.5 \\
Arrow of time & Asymptotic compactification per step & 7, 8 \\
Scale coincidence $R_{\rm eff}/\sigma\to 1$ & Compactification + Gaussian width & 8.2 \\
Second cascade constant $\operatorname{erf}(1/\sqrt{2})$ & Scale coincidence at $d\to\infty$ & 8.2 \\
Irreversibility is asymptotic & $R_{\rm eff}>0$; no sharp loss & 8.3 \\
Unitarity + apparent irreversibility & Global unitary; boundary retains info & 8.4 \\
Decoherence $D=e^{-\Delta x^2/4R_{\rm eff}^2}$ & Tracing fibre with slicing weight & 8.5 \\
Cascade coherence $93.9\%$ & Gram eigenvalue deficit $\varepsilon=6.1\%$ & 8.6 \\
$\hbar$ as quantum scale & Unit matching: $\hbar=(\pi/2)/\Delta t$ & 9 \\
Classical limit & Coarse time resolution $\Rightarrow$ no phase & 9.2 \\
Tensor product structure & Iterated slicing decomposition & 10 \\
Entanglement & Non-factorisable cascade states & 10.2 \\
Bell violation $S=2\sqrt{2}$ & Born rule on $S^7$ + CHSH bound & 10.3 \\
\bottomrule
\end{tabular}
\end{center}

Every row is derived from two inputs: (1) the cascade geometry of~\cite{paperI}, and
(2) the assumption that the observer has access to 4 dimensions. The compactification
results (Section~\ref{sec:compactification}) are proved in this paper from the Beta
function, building on the boundary dominance identity of~\cite{paperI}, Theorem~3.1.
The phase--obstruction lockstep (Corollary~\ref{cor:lockstep}) is derived from the
forced precession and sphere parity alone; the Clifford refinement to period~8 is
deferred to Part~III.

\section{What This Paper Does and Does Not Do}

\textbf{Does:}
\begin{itemize}
\item Shows that a 4D observer in the cascade geometry recovers the structural framework of
quantum mechanics.
\item Derives the Born rule $p = \cos^2\theta$ from the geometry of the spherical cap, the
slicing integrand, and Parseval's identity on $S^{d-1}$ (Theorem~\ref{thm:born}). No
probability axiom enters.
\item Proves the commutator of rank-1 projections:
$\|[P_{e_1},P_{e_2}]\| = \tfrac{1}{2}|\sin 2\theta|$
(Theorem~\ref{thm:commutator}), vanishing at $\theta\in\{0,\pi/2,\pi\}$ with maximum
at $\pi/4$.
\item Proves Bell inequality violation $S = 2\sqrt{2}$ on $S^7$
(Theorem~\ref{thm:bell}) from the Born rule applied to tensor product directions.
\item Proves that the precession angle $\alpha=\pi/2$ is forced by the cascade's
orthogonality axiom, eliminating all free geometric parameters
(Theorem~\ref{thm:precession}). The proof explicitly distinguishes the two applications of
the orthogonality axiom (one forcing $\sqrt{\pi}$ in the integral, one forcing $\alpha$)
and shows they are sequential applications of the same axiom, not a circular argument.
\item Derives the Gaussian character of wavefunctions from projection, with the elementary
mechanism identified: the slicing integrand $(1-x^2)^{d/2}$ is itself approximately Gaussian
(Lemma~\ref{lem:gaussian}).
\item Provides a geometric origin for complex amplitudes via forced axis precession.
\item Identifies each cascade slicing operation as one elementary unit of time
(dimensional collapse, \S\ref{subsec:time-id}): slicing \emph{is} dimensional
collapse and dimensional collapse \emph{is} time, yielding the arrow of time
and a natural lapse function. The geometric reading is asymptotic infall of
the observer's $S^3$-host shell toward the cover-sheet pseudo-horizon.
\item Closes the $U(1)$ vs $\mathbb{Z}_2$ gauge question
(Theorem~\ref{thm:U1-gauge-forced}): under the cascade hypothesis, the gauge
group on the projective state space is $U(1)$, the projective state space is
$\mathbb{C}P^{n-1}$ via the Hopf bundle, and the Born rule is the complex
$|\langle u, v\rangle_{\mathbb{C}}|^2$. Empirical commitment: the cascade is
falsifiable by a tripartite Bell experiment yielding real-amplitude
statistics (Renou \emph{et al.}~\cite{renou2021}).
\item Derives the discrete evolution equation (Theorem~\ref{thm:discrete-schrodinger}) and
shows the Schr\"{o}dinger equation arises as the effective description for a 4D observer
(Corollary~\ref{cor:schrodinger}), with $\mathcal{H}$ expressed as a Gamma function ratio.
The approximation is valid deep in the cascade; at $d=4$ the discrete propagator is primary.
\item Closes the cascade's \emph{kinematic} dynamics
(\S\ref{subsec:dynamics-closure}): the discrete propagator
(Theorem~\ref{thm:propagator}) plus Paper~VI's tower-growth
operator, under the period-1 tick reading
(forced cascade-internally --- period~2 excluded by the $U(1)$
gauge equivalence of Theorem~\ref{thm:U1-gauge-forced}, period~4
excluded by Paper~VI's inflation-window constraint), supplies the
cascade's full free-particle and cosmic-scale evolution. The
``missing Hamiltonian'' objection misreads where the cascade's
dynamical structure sits: $\mathcal{H}_{\rm eff}(d) = (1-N(d))/N(d)^2$
is c-number by construction, not a missing operator-valued
generator. The \emph{interaction} dynamics --- the cascade-native
gauge-coupled fermion action of Part~IVb~\cite{paper4b}
($\mathit{rem{:}fermion\text{-}gauge\text{-}coupling\text{-}proposal}$) ---
is structurally determined by six cascade-internal conditions, with
all originally-listed downstream items (Y-spectrum, $1/\alpha_{\rm em}$
screening, SU(2)$_L$ parity violation) now closed.
\item Proves the compactification radius $R_{\rm eff}=1/\sqrt{d+3}$ from the Beta function
(Theorem~\ref{thm:reff}) and develops its physical interpretation: each temporal step is an
asymptotic compactification. The arrow of time, the nature of measurement, and the
compatibility of unitarity with apparent irreversibility all receive concrete geometric
content from this identification.
\item Establishes the scale coincidence $R_{\rm eff}/\sigma\to 1$ (Theorem~\ref{thm:scale})
and derives the second universal cascade constant $\operatorname{erf}(1/\sqrt{2}) = 0.68269$
(Corollary~\ref{cor:erf}).
\item Shows that boundary dominance ($\Omega_{d-1}/V_d=d$) justifies the primacy of the
unit sphere as the observer's state space.
\item Derives the phase--obstruction lockstep (Corollary~\ref{cor:lockstep}): the
propagator phase is imaginary if and only if the sphere carries a hairy ball zero, both
controlled by the parity of $d$. The cascade-native classification has period~4; the
Clifford refinement to period~8 is noted and deferred to Part~III.
\item Derives the geometric decoherence rate (Theorem~\ref{thm:decoherence}): the
cascade's slicing traces out the fibre, producing Gaussian decoherence
$D = \exp(-\Delta x^2/4R_{\rm eff}^2)$ with decoherence length $2R_{\rm eff}$. The
structure matches environment-induced decoherence, with the integrated-out direction as the
environment and the compactification radius as the coherence length.
\end{itemize}

\textbf{Does not:}
\begin{itemize}
\item Derive why we observe 4 dimensions. This is an empirical input.
\item Derive the specific Hamiltonian of any physical system. The cascade gives the
framework but not the content.
\item Derive the value of $\hbar$. The geometric factor $\pi/2$ is derived; the physical
time increment $\Delta t$ is not.
\item Derive quantum field theory. The cascade gives single-particle QM; extension to fields
requires additional structure.
\item Use any result from physics. The only inputs are the cascade's geometry and the
observer's dimensionality.
\end{itemize}

\section{Open Questions}
\label{sec:open}

\begin{enumerate}
\item \textbf{Quantum field theory.} Whether the dimensions between successive thresholds
encode specific field content is open.
\item \textbf{Physical time increment.} The geometric value $\alpha=\pi/2$ is derived. The
remaining question is how the physical time increment $\Delta t$ is fixed, which would
determine $\hbar$.
\item \textbf{Lorentzian signature.} The cascade is real and Euclidean. The forced precession
$\alpha=\pi/2$ produces $J^2=-\mathrm{Id}$. Whether the identification $g_{tt}=(iN)^2=-N^2$
can be derived as a theorem is open.
\item \textbf{The measurement problem.} The cascade provides a natural account of measurement
as asymptotic compactification: the measured direction is suppressed to $R_{\rm eff}=
1/\sqrt{d+3}$, not sharply projected. Whether this smooth suppression---combined with the
boundary's retention of information ($\Omega/V=d$)---replaces the projection postulate is
open. The structural ingredients (global unitarity, apparent local irreversibility, boundary
encoding) are suggestive but not yet a derivation.
\item \textbf{Decoherence and the descent correction.}
Theorems~\ref{thm:decoherence} and~\ref{thm:cumulative-dec} derive the per-step and
cumulative decoherence: $\varepsilon = 6.11\%$ over the full 213-step descent, with the
cascade retaining $93.9\%$ coherence (Corollary~\ref{cor:cascade-coherence}). This
eigenvalue deficit is the same quantity that drives the descent corrections in the Part~0
Supplement. The remaining question: whether the cascade's geometric decoherence, applied to
specific physical systems, reproduces the quantitative predictions of environment-induced
decoherence~\cite{zurek}.
\item \textbf{$U(1)$ gauge upgrade: forced or permitted?}\label{oq:U1-gauge-upgrade}
\textsc{Resolved} by Theorem~\ref{thm:U1-gauge-forced} (\S7.5) and
Remark~\ref{rem:U1-closure-status}. The $U(1)$ upgrade is \emph{forced},
conditional on the cascade hypothesis: $\mathbb{Z}_2$ gauge predicts an
observable 2-cycle ray rotation per cascade step from a scalar
Hamiltonian, contradicting the standard-QM behaviour the hypothesis
requires; $U(1)$ gauge gives unobservable global phase, matching
standard QM. The closure uses only Part~II content
(Theorem~\ref{thm:complex}, Theorem~\ref{thm:discrete-schrodinger},
Theorem~\ref{thm:propagator}, Corollary~\ref{cor:schrodinger}, the
hypothesis of Section~1). Earlier ``both gauges cascade-consistent''
framing in Remarks~\ref{rem:U1-gauge-Born} and~\ref{rem:U1-gauge-J} is
superseded.
\item \textbf{Dim-2 Born-rule uniqueness via cascade compositionality.}\label{oq:dim2-born-uniqueness}
Remark~\ref{rem:U1-gauge-Born} establishes that under the now-forced
$U(1)$ gauge (Theorem~\ref{thm:U1-gauge-forced}), Cauchy additivity
on the $U(1)$-quotient simplex uniquely fixes the Born rule
$g(x) = x$ for $n \geq 3$. At $n = 2$ (qubit), single-system Cauchy
is under-determined exactly as in Gleason's theorem; closure was
deferred to ``compositionality with higher-dim cascade content
(the matter-content path tensor product of Part~IVa Section~4
reaches $n \geq 3$ for any composite system through gauge layers).''
This compositionality argument has not been formalised: a cascade-
internal proof that the tensor product of any cascade-derivable
$n = 2$ subsystem with cascade-internal higher-dim content reaches
the $n \geq 3$ regime, propagating Cauchy uniqueness back to $n = 2$,
remains open. Numerical verifier
\texttt{tools/verifiers/cascade\_path\_state\_born\_rule.py}
confirms the $n \geq 3$ uniqueness and the $n = 2$ gap; the gap is
narrowed by the closure of $oq$:U1-gauge-upgrade (which fixes the
gauge whose simplex is the right one for Cauchy) but not eliminated.
\end{enumerate}

\begin{thebibliography}{9}
\bibitem{paperI} RTAC, \emph{The Cascade Series --- Part 0: Scale Variance from
Orthogonality: How the Unit Ball Generates $10^{120}$ Orders of Magnitude}, March 2026.
\bibitem{ledoux} M.~Ledoux, \emph{The Concentration of Measure Phenomenon}, AMS
Mathematical Surveys and Monographs 89, 2001.
\bibitem{milman} V.~D.~Milman and G.~Schechtman, \emph{Asymptotic Theory of
Finite-Dimensional Normed Spaces}, Lecture Notes in Mathematics 1200, Springer, 1986.
\bibitem{fuchs} C.~A.~Fuchs, ``Quantum Mechanics as Quantum Information (and only a little
more),'' arXiv:quant-ph/0205039, 2002.
\bibitem{clifton} R.~Clifton, J.~Bub, and H.~Halvorson, ``Characterizing quantum theory in
terms of information-theoretic constraints,'' \emph{Foundations of Physics} \textbf{33},
1561--1591 (2003).
\bibitem{gleason} A.~M.~Gleason, ``Measures on the closed subspaces of a Hilbert space,''
\emph{Journal of Mathematics and Mechanics} \textbf{6}, 885--893 (1957).
\bibitem{zurek} W.~H.~Zurek, ``Quantum Darwinism,'' \emph{Nature Physics} \textbf{5},
181--188 (2009).
\bibitem{hardy} L.~Hardy, ``Quantum theory from five reasonable axioms,''
arXiv:quant-ph/0101012, 2001.
\bibitem{adm} R.~Arnowitt, S.~Deser, and C.~W.~Misner, ``The dynamics of general
relativity,'' in \emph{Gravitation: An Introduction to Current Research}, Wiley, 1962.
\bibitem{renou2021} M.-O.~Renou \emph{et al.}, ``Quantum theory based on real
numbers can be experimentally falsified,'' \emph{Nature} \textbf{600}, 625--629
(2021), arXiv:2101.10873.
\bibitem{paper4b} RTAC, \emph{The Cascade Series --- Part IVb: The Standard
Model from the Cascade --- Masses, Couplings, and Precision Predictions from
the Geometric-Topological Factorization}, March 2026.
\bibitem{paper6} RTAC, \emph{The Cascade Series --- Part VI: Tower Growth,
Inflation, and the Pre-Big-Bang from the Cascade}, April 2026.
\end{thebibliography}

\end{document}